\documentclass[11pt]{amsart}
\IfFileExists{lmodern.sty}%
  {\usepackage[T1]{fontenc}\usepackage{lmodern}\usepackage{microtype}}%
  {\usepackage[expansion=false]{microtype}}
\usepackage{amssymb,mathtools}
\usepackage{booktabs,array}
\usepackage{enumitem}
\emergencystretch=2.5em
\DeclareMathSizes{10}{10}{7}{6}
\DeclareMathSizes{9}{9}{6}{6}
\DeclareMathSizes{8}{8}{6}{6}
\usepackage[colorlinks=true,linkcolor=blue,citecolor=blue,urlcolor=blue]{hyperref}
\hypersetup{pdftitle={Irrationality measures for a zeta(3) + b pi\string^4 and zeta(3)/pi\string^4},pdfauthor={Sean Litsey}}

% ---------------------------------------------------------------------------
% Location of the Lean formalization.
\newcommand{\FormalizationURL}{https://github.com/slitsey/catalan-irrationality-measure}
\newcommand{\FormalizationLink}{\url{\FormalizationURL}}
% ---------------------------------------------------------------------------

\newtheorem{theorem}{Theorem}[section]
\newtheorem{lemma}[theorem]{Lemma}
\newtheorem{proposition}[theorem]{Proposition}
\newtheorem{corollary}[theorem]{Corollary}
\theoremstyle{definition}
\newtheorem{definition}[theorem]{Definition}
\theoremstyle{remark}
\newtheorem{remark}[theorem]{Remark}
\numberwithin{equation}{section}
\setcounter{tocdepth}{1}

\newcommand{\Z}{\mathbb Z}
\newcommand{\Q}{\mathbb Q}
\newcommand{\R}{\mathbb R}
\newcommand{\C}{\mathbb C}
\newcommand{\F}{\mathbb F}
\newcommand{\N}{\mathbb N}
\newcommand{\lean}[1]{\nolinkurl{#1}}
\newcommand{\vG}{v_p}
\newcommand{\cL}{\mathcal L}
\newcommand{\cO}{\mathcal O}
\newcommand{\cM}{\mathcal M}
\newcommand{\fr}[1]{\{#1\}}
\newcommand{\Lex}{\operatorname{E}}

\begin{document}

\title[Measures for $a\zeta(3)+b\pi^4$ and $\zeta(3)/\pi^4$]{Irrationality measures for $a\zeta(3)+b\pi^4$
and $\zeta(3)/\pi^4$}

\author{Sean Litsey}
\email{sean.litsey.math@gmail.com}

\date{November 17, 2026}

\subjclass[2020]{Primary 11J82; Secondary 11J72, 11M06, 15A15, 31A15, 68V20}

\keywords{Irrationality exponent, $\zeta(3)$, $\pi^4$, Hankel determinants, $p$-adic valuations, logarithmic
energy, formal verification, Lean}

\begin{abstract}
Let $\mu(\xi)$ denote the irrationality exponent of a real number $\xi$. We prove that
$\mu(a\zeta(3)+b\pi^4)\le27.49$ for all integers $a$ and $b$ with $b\ne0$, and that $\mu(\zeta(3)/\pi^4)\le35.35$. The
second bound extends to every quotient $(a\zeta(3)+b\pi^4)/(c\zeta(3)+d\pi^4)$ with rational $a,b,c,d$ and
$ad-bc\ne0$. Lai proved that $1$, $\zeta(3)$ and $\pi^4$ are linearly independent over $\Q$ with Hankel determinants
$J_n(a,b,X)$. We keep his determinants and prove the three estimates that a measure needs. First, the determinants
are small on a neighbourhood of the constant, uniformly in the direction $(a,b)$. The proof uses a positive measure
on a complex contour whose potential has a closed form. Second, the denominators are bounded prime by prime by an
exact law for the $p$-adic valuation. The integral of this law over the primes has a closed form. Third, $J_n$ does
not vanish at sizes $n$ tied to several families of primes. The main family consists of the primes $\ell$ with
$8n/3<\ell<3n-1$. For these primes a reduced Gram matrix is block triangular modulo $\ell$. Its half-integral block is
a Hankel determinant of Bernoulli-type moments, which we evaluate as a product of norms of Wilson polynomials. A
counting argument with the prime number theorem in progressions turns this nonvanishing into a measure. The best
published bound on the plane was $66.3$, and no measure was known for $\zeta(3)/\pi^4$. The proofs are formalized in
Lean~4.
\end{abstract}

\maketitle
\tableofcontents

%% =====================================================================
\section{Introduction}\label{sec:intro}

\subsection{The result}
The \emph{irrationality exponent} $\mu(\xi)$ of a real number $\xi$ is the infimum of the real numbers $\mu$ such
that $|\xi-p/q|>q^{-\mu}$ for all integers $p$ and all sufficiently large integers $q$. It does not change when
$\xi$ is multiplied by a nonzero rational number or translated by a rational number. A finite bound for $\mu(\xi)$
is an irrationality measure for $\xi$, and it shows that $\xi$ is irrational. We write
\[
 \zeta(3)=\sum_{k\ge1}\frac1{k^3}=1.20205\,69031\ldots,\qquad \zeta(4)=\sum_{k\ge1}\frac1{k^4}=\frac{\pi^4}{90}.
\]

\begin{theorem}\label{thm:main}
\begin{enumerate}[label=(\alph*)]
\item Let $a$ and $b$ be integers with $b\ne0$. There is an integer $q_0$ such that
\[
 \Bigl|a\zeta(3)+b\pi^4-\frac pq\Bigr|>q^{-27.49}
\]
for all integers $p$ and all integers $q\ge q_0$. In particular $\mu(a\zeta(3)+b\pi^4)\le27.49$.
\item There is an integer $q_0$ such that
\[
 \Bigl|\frac{\zeta(3)}{\pi^4}-\frac pq\Bigr|>q^{-35.35}
\]
for all integers $p$ and all integers $q\ge q_0$. In particular $\mu(\zeta(3)/\pi^4)\le35.35$.
\end{enumerate}
\end{theorem}

Since $\mu$ is invariant under multiplication by nonzero rationals, part (a) bounds the exponent of every number
$r\zeta(3)+r'\pi^4$ with $r,r'\in\Q$ and $r'\ne0$. The integer $q_0$ of part (a) depends on $a$ and $b$.
The exponent is also invariant under rational M\"obius maps, so part (b) extends to every quotient of two
independent points of the plane.

\begin{corollary}\label{cor:mobius}
Let $a$, $b$, $c$ and $d$ be integers with $ad-bc\ne0$. There is an integer $q_0$ such that
\[
 \Bigl|\frac{a\zeta(3)+b\pi^4}{c\zeta(3)+d\pi^4}-\frac pq\Bigr|>q^{-35.35}
\]
for all integers $p$ and all integers $q\ge q_0$. In particular
$\mu\bigl((a\zeta(3)+b\pi^4)/(c\zeta(3)+d\pi^4)\bigr)\le35.35$.
\end{corollary}

\begin{proof}
Put $y^*=\zeta(3)/\zeta(4)$. Since $\pi^4=90\zeta(4)$, the number is $(ay^*+90b)/(cy^*+90d)$, and
$90(ad-bc)\ne0$. Section~\ref{sec:proof} gives $|y^*-p/q|>q^{-35.3495}$ for all integers $p$ and all large $q$.
Lemma~\ref{lem:mobius} applies with $\mu'=35.3495$ and $\mu=35.35$.
\end{proof}

The corollary also holds for rational $a$, $b$, $c$ and $d$. Multiplying them by a common denominator changes
neither the number nor the condition $ad-bc\ne0$. The proof gives the exponents $27.4849$ and $35.3492$, with the
constants of Table~\ref{tab:constants}.
Theorem~\ref{thm:main} and Corollary~\ref{cor:mobius} have been verified in the Lean~4 proof assistant, together
with every theorem, proposition and lemma used in their proofs. The formal proof has no unproved steps and uses only the standard axioms.
Section~\ref{sec:lean} describes it. The remarks are not part of the formal proof. The formalization is available at
\FormalizationLink. The use of AI in this work is disclosed at the end of the paper.

\subsection{Background}\label{ssec:background}
\emph{The numbers $\zeta(3)$ and $\pi^4$.} Euler showed that $\zeta(4)=\pi^4/90$, and $\pi^4$ is
transcendental. Ap\'ery proved that $\zeta(3)$ is irrational \cite{Apery79}, and Beukers gave his integral proof
\cite{Beukers79}. Rhin and Viola introduced a method of permutation groups for $\zeta(2)$ \cite{RhinViola96}, and
with it they proved $\mu(\zeta(3))\le5.513891$ \cite{RhinViola01}. For $\pi^4$
the best bound is due to Marcovecchio and Zudilin, $\mu(\pi^4)=\mu(\zeta(4))\le12.51085940\ldots$ \cite{MarcovecchioZudilin20}.
For the other odd values the classical results concern families. Rivoal \cite{Rivoal00} and Ball and Rivoal
\cite{BallRivoal01} proved that infinitely many of the numbers $\zeta(2k+1)$ are irrational. Zudilin proved that
one of $\zeta(5)$, $\zeta(7)$, $\zeta(9)$ and $\zeta(11)$ is irrational \cite{Zudilin01,Zudilin04}. Fischler, Sprang
and Zudilin raised the number of irrational values among $\zeta(3),\dots,\zeta(s)$ \cite{FSZ19}. Fischler proved that
many of the odd zeta values are linearly independent over $\Q$ \cite{Fischler26}.

\emph{The plane spanned by $\zeta(3)$ and $\pi^4$.} Whether $1$, $\zeta(3)$ and $\pi^4$ are linearly
independent over $\Q$ was a classical open question. Dauguet and Zudilin studied simultaneous rational
approximations to $\zeta(2)$ and $\zeta(3)$ of the hypergeometric type, and they noted that the independence of $1$,
$\zeta(2)$ and $\zeta(3)$ remained open \cite{DauguetZudilin14}. Fauzan proved that $\zeta(5)$ is irrational with
rationally normalized Hankel determinants \cite{Fauzan2026}. Lai combined Fauzan's determinant method with the
simultaneous approximations of Zudilin and of Dauguet and Zudilin, and proved that $1$, $\zeta(3)$ and $\pi^4$ are
linearly independent \cite{Lai2026}. As a corollary $\zeta(3)/\pi^4$ is irrational. Leibtag, Barkan, Kalisch,
Weinbaum, Shalyt and Kaminer connected the classical constructions through conservative matrix fields
\cite{LBKWSK26}. They gave another proof of this independence and proved $\mu(A\zeta(3)+B\pi^4)<66.3$ for every
$(A,B)\in\Q^2\setminus\{(0,0)\}$. They also proved $\mu(\zeta(2))<4.62675$ and $\mu(\zeta(3))<5.4197$. Fischler, Lai,
Sprang and Zudilin proved that many more odd zeta values are irrational \cite{FLSZ26}. As far as we know, $66.3$ is
the best published measure on the plane, and no irrationality measure for $\zeta(3)/\pi^4$ was known. Part (a)
improves the bound to $27.49$ for $b\ne0$. Part (b) is, as far as we know, the first measure for $\zeta(3)/\pi^4$. On the line $b=0$ the number is a rational multiple of $\zeta(3)$, whose
measure is much smaller. On the line $a=0$ the bound of Marcovecchio and Zudilin is smaller than ours.

\emph{Recent results.} Since September 2026 several new results have appeared in this area, some obtained with AI
systems. In October 2026 OpenAI proved that Catalan's constant is irrational \cite{OpenAI2026} and that $\mu(\pi)=2$
\cite{OpenAIPi2026}, both with Lean formalizations. Liu, Jiang and Zhang proved $\mu(\log r)=2$ for rational $r>0$,
$r\ne1$ \cite{LiuJiangZhang26}. Giraudin improved the bounds for $\mu(\zeta(2))$ and $\mu(\zeta(3))$, and he gave
a bound for $\mu(\zeta(5))$ \cite{Giraudin26,Giraudin26z3,Giraudin26z5}. Q.~Tang proved the linear independence of $1$,
$\zeta(2)$ and $\zeta(3)$ in Lean \cite{Tang26}.

\emph{The method.} The measure method we use, a family of integer polynomials whose values near the constant are
small and nonzero, is a one-variable form of the criteria of Nesterenko \cite{Nesterenko85} used by Ball and Rivoal
\cite{BallRivoal01}. The arithmetic of the coefficients is handled prime by prime, as in the work of Hata
\cite{Hata93}, Rhin and Viola \cite{RhinViola96,RhinViola01} and Zudilin \cite{Zudilin04}. Determinants of growing
size go back to Zudilin's determinantal approach \cite{Zudilin}. Their size is controlled by Andr\'eief's identity
\cite{Forrester2018Andreief} and a logarithmic energy problem with an external field \cite{SaffTotik}.

\subsection{The companion papers}
This paper applies one method to a fourth set of constants. The method consists of Hankel or L\"owner determinants,
denominators bounded prime by prime, decay proved through an energy problem, and a formal proof of every step. The
first companion paper \cite{Litsey2026} refines OpenAI's determinant for Catalan's constant. The second
\cite{Litsey2026b} uses L\"owner determinants of digamma combinations and proves measures on three planes of numbers.
Its third plane is the plane of this paper, since $\pi^4=90\zeta(4)$. The third \cite{Litsey2026c} proves a measure
for $\zeta(5)$ with Fauzan's determinants. Its decay argument, a comparison measure with closed-form potentials checked
on cells, is the closest relative of Section~\ref{sec:decay}. We use the definitions of $\mu$, the measure criterion
and the potential-theoretic language of these papers. The construction here is Lai's, with his notation. From Section~\ref{ssec:require} on we write the
points of the plane as $a\zeta(3)+b\zeta(4)$, as Lai does. Since $\pi^4=90\zeta(4)$, the number $a\zeta(3)+b\pi^4$ of
Theorem~\ref{thm:main}(a) is $a\zeta(3)+90b\zeta(4)$.

\subsection{What a measure requires}\label{ssec:require}
Lai's proof uses three facts about his determinants $J_n=J_n(a,b,a\zeta(3)+b\zeta(4))$. A rational normalizing factor
$m_n$ with $\log m_n\le(30.83+o(1))n^2$ makes $m_nJ_n(a,b,X)$ a polynomial with integer coefficients. The decay
$|J_n|\le e^{-31.61n^2}$ holds at the constant. And $J_n\ne0$ for $n=2(\ell-1)/9$ at a prime $\ell\equiv1\pmod9$ with
$\ell>|b|$. If $a\zeta(3)+b\zeta(4)$ were rational, a nonzero integer would be smaller than $1$.

A measure needs more. Let $p/q$ be a rational close to the constant. Then $q^{3n}m_nJ_n(a,b,p/q)$ is an integer, and
three estimates are needed at sizes $n$ of order $\log q$.
\begin{enumerate}[label=(\roman*)]
\item \emph{Denominators.} The bound $\log m_n\le(\Lambda+o(1))n^2$.
\item \emph{Decay on a neighbourhood.} $|J_n(a,b,X)|\le e^{-(D-o(1))n^2}$ for every real $X$ within
 $e^{-(\rho+o(1))n}$ of the constant, not only at the constant.
\item \emph{Nonvanishing near the point.} $J_n(a,b,p/q)\ne0$ for one size $n$ in a range of length
 $O(\log q)$ above the size where (i) and (ii) beat $q^{3n}$.
\end{enumerate}
Lai's nonvanishing holds only at sizes $n\gg|b|$. For the ratio $\zeta(3)/\pi^4$ the rational $p/q$ enters through
the direction: $q^{3n}J_n(1,-p/q,0)=J_n(q,-p,0)$, so the coefficient $b=-p$ grows with $q$. So the decay must be
uniform in the direction $(a,b)$, and nonvanishing must hold at sizes of order $\log q$ for coefficients of size $q$.
The exponent that results is $\rho(3/(D-\Lambda)+\beta)$, where $\beta$ measures the sparsity of the sizes at which
nonvanishing is proved.

\subsection{New ideas}
We keep Lai's determinants unchanged. The proof rests on the following ideas.
\begin{enumerate}[leftmargin=*]
\item \emph{Decay on a complex contour by weak duality.} Lai writes the functional as an integral over a vertical
 line and bounds $J_n$ by Heine's formula with the modulus of a complex weight. We prove a general upper bound of
 energy type for this situation. A positive measure $\sigma$ of mass $3$ in the complex plane and an upper bound
 $W^*$ for $V-2U^\sigma$ on an admissible contour give $|J_n|\le e^{-(D-o(1))n^2}$ with $D=-\kappa-I(\sigma)-3W^*$.
 The perturbation of $X$ adds atoms at the negative integers. They are controlled by the potential of $\sigma$ on
 the negative axis, and this gives the neighbourhood exponent $\rho$. The bound is uniform in the direction $(a,b)$
 (Theorem~\ref{thm:dual}). We apply it to a contour with two straight arms and a measure on two segments whose
 density against the arcsine law is an explicit square. Its potentials have closed forms in real arithmetic, and the
 inequality is checked on $40$ cells (Theorem~\ref{thm:decay}).
\item \emph{An exact law for the denominators.} For every prime $\sqrt{8n}<p<3n$ we improve Lai's bound for
 $v_p(J_n)$ by two explicit integers $x_1$ and $x_2$, determined by the ties among Lai's increments
 (Theorems~\ref{thm:L1} and~\ref{thm:L2}). The first gain is a rank bound for a reduced Gram form modulo $p$. Its weights
 are values of one polynomial on residue classes, so a power sum over $\F_p$ forces a kernel. The second gain comes
 from a basis whose roots are shifted by multiples of $p$, so that the leading local forms are orthogonal. The
 remaining terms cancel modulo $p^2$, because a Faulhaber sum and the harmonic weights of the unit parts have
 opposite totals. For $3n<p\le6n$ similar bases give Lai's case bounds for the polynomial in three variables and an
 extra $7n-2p$ in the range $3n<p\le7n/2$ (Theorems~\ref{thm:O} and~\ref{thm:B7}).
\item \emph{The integral of the law in closed form.} On $p\le3n$ the law is piecewise affine in $n/p$ on $22$
 pieces. Its integral against the prime distribution telescopes to a sum over $23$ nodes, and Gauss's digamma
 theorem gives a closed form in logarithms (Theorem~\ref{thm:C01}).
\item \emph{Nonvanishing in a band of primes.} For every prime $\ell$ with $8n/3<\ell<3n-1$, outside a
 finite set, we prove $J_n(a,b,X)\ne0$ at every point with $a\in\Z_{(\ell)}$, $b$ an $\ell$-adic unit and
 $\ell^3X\in\Z_{(\ell)}$ (Theorem~\ref{thm:B83}). So a prime is good for every size $n$ in an interval, and not for one
 size only. The proof shows that a reduced Gram matrix is nonsingular modulo $\ell$. This matrix is block triangular. Its integral blocks are small matrices of rational
 constants with explicit determinants (Lemmas~\ref{lem:U83}, \ref{lem:blockA} and~\ref{lem:blockC}). Its half-integral block
 is, modulo $\ell$, the Hankel matrix of the functional $x^e\mapsto eB_{e-1}$ against the cube of a falling factorial
 (Lemma~\ref{lem:G3}). We obtain this by Faulhaber's formula and four sums of harmonic numbers, and not by analysis.
\item \emph{The Hankel determinant in closed form.} The determinant of the half-integral block is an explicit
 product of factorial ratios, the norms of two families of Wilson polynomials, and all its prime factors are
 smaller than $\ell$ (Lemma~\ref{lem:W}). We prove it by finite algebra: a Pearson equation for the functional, a
 Beta-type evaluation, Gram--Schmidt for the orthogonal polynomials and the Pfaff--Saalschütz summation. No integral
 is used.
\item \emph{Nonvanishing at prime-tied sizes.} We extend Lai's auxiliary-prime argument to every rational $X$ whose
 denominator is not divisible by $\ell^4$, and to the ratio direction (Theorem~\ref{thm:NV}). We find a second family
 of primes with exact nonvanishing, $3\ell=7n+1$ and $3\ell=7n+2$ (Theorem~\ref{thm:NV7}). There the reduction
 modulo $\ell$ leads to a Hankel determinant of a weight on an interval of residues. By duality for complementary
 subsets its value is a determinant of size $1$ or $2$ of power sums of a polynomial over $\F_\ell$, equal to
 $-1/49$ or $-1$.
\item \emph{A criterion with bands.} Suppose that nonvanishing is proved for families of primes in residue classes,
 and that each prime is good for the sizes $n$ in an interval. The prime number theorem in progressions and an
 inclusion--exclusion count over the classes turn this into a measure (Theorem~\ref{thm:criterion}). A prime in several families is counted once.
 The band of item~4 covers every prime between $8n_0/3$ and $3N$ for a window of sizes $n_0\le n\le N$.
\end{enumerate}

\begin{table}[ht]
\centering
\renewcommand{\arraystretch}{1.2}
\setlength{\tabcolsep}{4pt}
\begin{tabular}{@{}lccc>{\raggedright\arraybackslash}p{4.4cm}@{}}
\toprule
 & $\Lambda$ & $D$ & $\rho$ & nonvanishing\\
\midrule
Lai \cite{Lai2026} & $30.83$ & $31.61$ & -- & $n=2(\ell-1)/9$, $\ell>|b|$\\
This paper & $28.68723663$ & $33.4401$ & $42.94856306$ & band $\frac83n<\ell<3n-1$, $\ell\equiv1\ (18)$, $\ell\equiv3,5\ (7)$\\
\bottomrule
\end{tabular}
\caption{The constants per $n^2$ (and per $n$ for $\rho$). Lai's decay holds at the constant only. The exponent is
$\rho(3/(D-\Lambda)+\beta)$ with $\beta=7/800$ for the plane and $\beta=9593/50000$ for the ratio.}\label{tab:constants}
\end{table}

\subsection{Organization}
Section~\ref{sec:construction} recalls Lai's construction. Section~\ref{sec:criterion} proves the measure criterion.
Section~\ref{sec:decay} proves the decay. Sections~\ref{sec:local}, \ref{sec:second} and~\ref{sec:sum} bound the
normalizing factor: the local bounds, the second level of the law, and the prime sum. Section~\ref{sec:nv} proves the
nonvanishing, in Lai's sizes, on two lines, and in the band of primes with its Hankel determinant. Section~\ref{sec:proof} completes the proof, and Section~\ref{sec:remarks} compares the constants with
numerical values. Section~\ref{sec:lean} describes the formalization. The disclosure of AI use follows it.

\subsection{Notation}
For a prime $p$ we write $v_p$ for the Gauss valuation on $\Q(a,b,X)$: the $p$-adic valuation extended to
polynomials by the minimum over the coefficients, with $v_p(a)=v_p(b)=v_p(X)=0$. We write $\Z_{(p)}$ for the
rationals with $v_p\ge0$, $H_m=\sum_{i\le m}1/i$ and $H_m(j)=\sum_{i\le m}i^{-j}$. We write $\fr y=y-\lfloor y\rfloor$
and $y_+=\max(y,0)$. Chebyshev's function in a progression is $\theta(x;M,r)=\sum_{\ell\le x,\ \ell\equiv r\ (M)}\log\ell$,
summed over primes $\ell$. The logarithmic potential and energy of a positive measure $\sigma$ on $\C$ are
\[
 U^\sigma(s)=\int\log\frac1{|s-u|}\,d\sigma(u),\qquad I(\sigma)=\iint\log\frac1{|s-u|}\,d\sigma(s)\,d\sigma(u).
\]

%% =====================================================================
\section{Lai's construction}\label{sec:construction}

We recall the construction of \cite[\S2]{Lai2026}, with the variable $X$ in place of the constant. Throughout $n\ge1$.
Write $(x)_m=x(x+1)\cdots(x+m-1)$ and put
\[
 R_n(t)=\frac{(6n)!^2}{(2n)!^3}\cdot\frac{(t-2n)_{2n}^3}{(t)_{6n+1}^2}.
\]
It has zeros of order $3$ at $t=1,\dots,2n$, double poles at $t=0,-1,\dots,-6n$, and $R_n(t)=O(t^{-6n-2})$. For
polynomials $P$ of degree less than $6n$ Lai defines
\begin{equation}\label{eq:Lseries}
 L_n(P)=-3a\sum_{m\ge1}(R_nP)'(m)+b\sum_{m\ge1}(R_nP)''(m),
\end{equation}
and
\[
 J_n=\frac{\det\bigl[L_n(t^{i+j})\bigr]_{0\le i,j<3n}}{\prod_{i<3n}i!^2}.
\]
Put
\[
 c_k=\binom{6n}k^2\binom{2n+k}k^3,\qquad d_k=5H_k-3H_{2n+k}-2H_{6n-k}\qquad(0\le k\le6n).
\]

\begin{proposition}[Lai]\label{prop:finite}
For $P$ of degree less than $6n$,
\begin{equation}\label{eq:finite}
 L_n(P)=\sum_{k=0}^{6n}c_k\bigl(U_kP'(-k)+V_kP(-k)\bigr),
\end{equation}
where $U_k=-3aH_k(2)-2bH_k(3)$ and $V_k=U_kd_k+6\bigl(X-aH_k(3)-bH_k(4)\bigr)$ at $X=a\zeta(3)+b\zeta(4)$.
\end{proposition}

\begin{proof}
The partial fraction expansion of $R_nP$ has the form
$\sum_k\bigl(A_{2,k}(t+k)^{-2}+A_{1,k}(t+k)^{-1}\bigr)$, since $\deg(R_nP)\le-2$. Here
$A_{2,k}=c_kP(-k)$, because $(t+k)^2R_n$ takes the value $c_k$ at $t=-k$. The logarithmic derivative of
$(t+k)^2R_n$ at $-k$ is $-3(H_{2n+k}-H_k)-2(H_{6n-k}-H_k)=d_k$, so $A_{1,k}=c_k(P'(-k)+d_kP(-k))$. Also
$\sum_kA_{1,k}=0$. Now $\sum_{m\ge1}(m+k)^{-s}=\zeta(s)-H_k(s)$. Differentiating the expansion once and twice and
summing over $m$, the terms in $\zeta(2)$ cancel because $\sum_kA_{1,k}=0$. The term $2b\zeta(3)\sum_kA_{1,k}$
cancels for the same reason. What remains is $6(a\zeta(3)+b\zeta(4))\sum_kA_{2,k}$ plus the stated rational part.
\end{proof}

We take \eqref{eq:finite} with $X$ an indeterminate as the definition of a $\Q[a,b,X]$-valued functional, and we write
$J_n(a,b,X)$ for the determinant above with this functional. The entries are linear forms in $(a,b,X)$, so
$J_n(a,b,X)\in\Q[a,b,X]$ is homogeneous of degree $3n$:
\begin{equation}\label{eq:homog}
 J_n(\lambda a,\lambda b,\lambda X)=\lambda^{3n}J_n(a,b,X).
\end{equation}
At $X=a\zeta(3)+b\zeta(4)$ it is Lai's $J_n$. The two families of this paper are
\[
 Q_n(X)=J_n(a,b,X)\quad\text{for fixed }a,b\in\Z,\qquad f_n(y)=J_n(1,-y,0).
\]
Both have degree at most $3n$. By \eqref{eq:homog},
\begin{equation}\label{eq:scaled}
 q^{3n}Q_n(p/q)=J_n(qa,qb,p),\qquad q^{3n}f_n(p/q)=J_n(q,-p,0).
\end{equation}
The constant of $Q_n$ is $a\zeta(3)+b\zeta(4)$. The constant of $f_n$ is $y^*=\zeta(3)/\zeta(4)=1.11062\ldots$, and
$\zeta(3)/\pi^4=y^*/90$. Indeed $f_n(y)=J_n(1,-y,X_y+t)$ with $X_y=\zeta(3)-y\zeta(4)$ and $t=-X_y=\zeta(4)(y-y^*)$.
So $f_n$ is small near $y^*$ when $J_n$ is small near the constant, uniformly in the direction.

%% =====================================================================
\section{The measure criterion}\label{sec:criterion}

The criterion of the companion papers \cite{Litsey2026,Litsey2026c} assumes nonvanishing at every large size. Here
nonvanishing is known only at the sizes $n$ tied to primes $\ell$ in some residue classes. For some families each
prime is tied to one size, and these sizes have density zero. For others each prime is tied to an interval of sizes
of length proportional to $\ell$. In both cases the primes carry enough logarithmic mass. If $P_n(p/q)=0$ at every tied
size in a window, then a power of every tied prime divides $q$, and the product of these powers exceeds $q$. The
criterion makes this precise.

Fix a modulus $M$. A \emph{tied family} consists of a set $S$ of unit classes modulo $M$, slopes
$0<c\le c'$, an integer $C\ge0$, an exponent $e\ge1$ and a finite set $E$ of primes. For every prime $\ell\notin E$
with $\ell\bmod M\in S$ it has a set $R(\ell)$ of sizes $n\ge1$, with the following property. For every real
$y\ge1$ with $cy\le\ell\le c'y$ the set $R(\ell)$ contains an integer $n$ with $|n-y|\le C$. A \emph{line} has $c=c'$ and
$R(\ell)=\{n(\ell)\}$ with $\ell=c\,n(\ell)+d$. A \emph{band} has $c<c'$, and $R(\ell)$ is the set of all $n$ with
$cn<\ell<c'n-d$. For a list of tied families $i=1,\dots,r$ and a nonempty subset $T$ of
indices put
\[
 \delta_T=\frac{\#\{s\in(\Z/M)^\times:s\in S_i\text{ for all }i\in T\}}{\varphi(M)},
\]
\[
 \lambda_T(u,w)=\Bigl(w\min_{i\in T}c'_i-u\max_{i\in T}c_i\Bigr)_+,
\]
and
\[
 \cM(u,w)=\sum_{\emptyset\ne T\subseteq\{1,\dots,r\}}(-1)^{|T|+1}\Bigl(\min_{i\in T}e_i\Bigr)\delta_T\,\lambda_T(u,w).
\]

\begin{theorem}[measure criterion]\label{thm:criterion}
Let $\xi$ be real, $\kappa\ge1$ an integer, and $\Lambda<D$, $\rho>0$, $\beta>0$, $K>0$ real numbers. Let $(P_n)_{n\ge1}$
be polynomials with rational coefficients and $\deg P_n\le\kappa n$. Suppose that for every $\varepsilon>0$ and all
large $n$ the following hold.
\begin{enumerate}[label=(\alph*)]
\item There is a positive rational number $d_n$ with $\log d_n\le(\Lambda+\varepsilon)n^2$ and $q^{\kappa n}d_nP_n(p/q)\in\Z$
 for all integers $p$ and $q>0$.
\item $|P_n(\xi+t)|\le e^{-(D-\varepsilon)n^2}$ for every real $t$ with $|t|\le e^{-(\rho+\varepsilon)n}$.
\end{enumerate}
Let $W(p,q)$ be positive integers with $W(p,q)\le Kq$ whenever $|\xi-p/q|<1$. Let tied families $i=1,\dots,r$ be
given, each with a relation $\mathrm{Bad}_i(\ell,p,q)$ such that:
\begin{enumerate}[label=(\alph*),resume]
\item for every prime $\ell\notin E_i$ with $\ell\bmod M\in S_i$, every $n\in R_i(\ell)$, and all $p$ and $q>0$, either
 $P_n(p/q)\ne0$ or $\mathrm{Bad}_i(\ell,p,q)$ holds, and $\mathrm{Bad}_i(\ell,p,q)$ implies
 $\ell^{e_i}\mid W(p,q)$;
\item $\cM(u_0,u_0+\beta)>1$, where $u_0=\kappa/(D-\Lambda)$.
\end{enumerate}
Then for every $\mu>\rho(\kappa/(D-\Lambda)+\beta)$ there is $q_0$ such that $|\xi-p/q|>q^{-\mu}$ for all integers $p$
and all integers $q\ge q_0$.
\end{theorem}

The fraction $p/q$ need not be in lowest terms, and no irrationality of $\xi$ is assumed. Condition (a) holds for
every representation of the fraction, and so does (c).

\begin{proof}
Put $\eta=D-\Lambda$ and fix $\mu>\rho(\kappa/\eta+\beta)$. Choose $\varepsilon>0$ so small that
$\cM(u_\varepsilon,w_\varepsilon)>1+\varepsilon$ and $\rho_\varepsilon(u_\varepsilon+\beta+\varepsilon)<\mu$, where
$u_\varepsilon=(\kappa+\varepsilon)/(\eta-2\varepsilon)+\varepsilon$, $w_\varepsilon=u_0+\beta-\varepsilon$ and
$\rho_\varepsilon=\rho+\varepsilon$. This is possible because $\cM$ is continuous: it is a finite sum of functions
$(w\gamma-u\gamma')_+$.

Let $q$ be large and suppose $|\xi-p/q|\le q^{-\mu}$. Put $L=\log q$,
\[
 n_0=\frac{(\kappa+\varepsilon)L+C}{\eta-2\varepsilon},\qquad N=n_0+(\beta+\varepsilon)L,
\]
with a constant $C$ chosen below. Then $\rho_\varepsilon N\le\mu L$ for large $q$, so $|\xi-p/q|\le
e^{-\rho_\varepsilon n}$ for every $n\le N$.

\emph{Step 1: no tied size in the window is nonzero.} Let $n\in[n_0,N]$ be large and suppose $P_n(p/q)\ne0$. By (a)
the number $z=q^{\kappa n}d_nP_n(p/q)$ is a nonzero integer, so $|z|\ge1$. By (b) with $t=p/q-\xi$,
$|z|\le q^{\kappa n}e^{(\Lambda+\varepsilon)n^2-(D-\varepsilon)n^2}$. Hence $(\eta-2\varepsilon)n\le\kappa L$. This
contradicts $n\ge n_0$ when $C>0$.

\emph{Step 2: one divisibility.} Let $U_i$ be the set of primes $\ell\notin E_i$ with $\ell\bmod M\in S_i$ for which
$R_i(\ell)$ meets $[n_0,N]$, and $U=\bigcup_iU_i$. For $\ell\in U$ let $e(\ell)$ be the largest $e_i$ with $\ell\in U_i$. By
Step 1 and (c), $\ell^{e_i}\mid W(p,q)$ for every $i$ with $\ell\in U_i$, so $\ell^{e(\ell)}\mid W(p,q)$. The primes are
distinct, so
\begin{equation}\label{eq:budget}
 \sum_{\ell\in U}e(\ell)\log\ell\le\log W(p,q)\le L+\log K.
\end{equation}

\emph{Step 3: the prime count.} Let $\ell\notin E_i$ be a prime with $\ell\bmod M\in S_i$ and
$c_i(n_0+C_i)<\ell\le c_i'(N-C_i)$. Then the interval $[\ell/c_i',\ell/c_i]\cap[n_0+C_i,N-C_i]$ is not empty. Take
$y$ in it. By the property of the family, $R_i(\ell)$ contains an integer $n$ with $|n-y|\le C_i$, and this $n$ lies in
$[n_0,N]$. So $\ell\in U_i$. For large $q$ the set $U_i$ therefore contains the primes of the class with $\ell/L$ in the interval
$J_i=(c_iu_\varepsilon,c_i'w_\varepsilon]$. For a fixed class $s$ the finitely many endpoints of the $J_i$ with
$s\in S_i$ cut $(0,\infty)$ into intervals. On each the set of $i$ with $J_i$ containing it is constant, and so is the
largest $e_i$. The prime number theorem in arithmetic progressions,
$\theta(x;M,s)=x/\varphi(M)+o(x)$ for every unit class $s$, counts the primes of each class in each interval. Removing
the finite sets $E_i$ costs $O(1)$. Summing,
\[
 \sum_{\ell\in U}e(\ell)\log\ell\ge\Bigl(\frac1{\varphi(M)}\sum_s\int_0^\infty\max_{i:\,s\in S_i}e_i\mathbf 1_{J_i}(x)\,dx
 -\frac\varepsilon2\Bigr)L
\]
for large $q$. For nonnegative numbers $x_i$, $\max_ix_i=\sum_{T\ne\emptyset}(-1)^{|T|+1}\min_{i\in T}x_i$. With
$x_i=e_i\mathbf 1_{J_i}(x)$, the minimum over $T$ is $(\min_Te_i)\mathbf 1_{\bigcap_TJ_i}(x)$, and
$\bigcap_TJ_i$ has length $\lambda_T(u_\varepsilon,w_\varepsilon)$. Averaging over $s$ gives the integral
$\cM(u_\varepsilon,w_\varepsilon)>1+\varepsilon$. So the left side of \eqref{eq:budget} is at least
$(1+\varepsilon/2)L$ for large $q$, which contradicts \eqref{eq:budget}.

Hence $|\xi-p/q|>q^{-\mu}$ for all large $q$. The constant $C$ absorbs $\log K$ and the finitely many small $n$.
\end{proof}

When the windows of different families do not overlap, $\cM(u,u+\beta)=\sum_ie_i\delta_i\bigl((c'_i-c_i)u+c'_i\beta\bigr)$. The inequality
(d) then says that the tied primes in a window of length $\beta L$ carry more than $L$ of logarithmic mass.

The second part of Theorem~\ref{thm:main} is stated for $\zeta(3)/\pi^4$. We apply the criterion to $y^*$ and use the
following step.

\begin{lemma}\label{lem:scale}
Let $x$ be real, $r\ne0$ rational and $\mu'<\mu$. Suppose that $|x-p/q|>q^{-\mu'}$ for all integers $p$ and all large
$q$. Then $|rx-p/q|>q^{-\mu}$ for all integers $p$ and all large $q$.
\end{lemma}

\begin{proof}
Write $r=u/v$ with integers $u\ne0$ and $v>0$. Then $|rx-p/q|=(|u|/v)\,|x-vp/(uq)|$, and $vp/(uq)$ is a fraction
with denominator $|u|q$. For large $q$ this is at least $(|u|/v)(|u|q)^{-\mu'}$, which exceeds $q^{-\mu}$ for large
$q$.
\end{proof}

Corollary~\ref{cor:mobius} uses the same step for integer M\"obius maps.

\begin{lemma}\label{lem:mobius}
Let $x$ be real, let $1\le\mu'<\mu$, and let $a$, $b$, $c$ and $d$ be integers with $ad-bc\ne0$. Suppose that
$|x-p/q|>q^{-\mu'}$ for all integers $p$ and all large $q$. Then $y=(ax+b)/(cx+d)$ satisfies $|y-p/q|>q^{-\mu}$ for
all integers $p$ and all large $q$.
\end{lemma}

\begin{proof}
The hypothesis makes $x$ irrational. Since $(c,d)\ne(0,0)$, the number $cx+d$ is not zero. Choose $q_1$ such that
the hypothesis holds for all $q\ge q_1$, and put $K=q_1$. For integers $Q\ne0$ and $P$ the fraction $P/Q$ has the
form $p'/(K|Q|)$, so
\begin{equation}\label{eq:allden}
 |Qx-P|=|Q|\,\Bigl|x-\frac PQ\Bigr|>|Q|\,(K|Q|)^{-\mu'}.
\end{equation}
Given $p$ and $q\ge1$, put $P=dp-bq$ and $Q=aq-cp$. Then $Qx-P=q(ax+b)-p(cx+d)$, so
\[
 y-\frac pq=\frac{Qx-P}{(cx+d)q}.
\]
The pair $(P,Q)$ is the image of $(p,q)\ne(0,0)$ under a matrix of determinant $ad-bc\ne0$, so $(P,Q)\ne(0,0)$.
If $Q=0$ then $|P|\ge1$ and $|y-p/q|\ge1/(|cx+d|q)$, which exceeds $q^{-\mu}$ for large $q$ since $\mu>1$. Let
$Q\ne0$. We may assume $|y-p/q|<1$. Then $|p|/q<|y|+1$, so $|Q|\le Cq$ with $C=|a|+|c|(|y|+1)$. Since $\mu'\ge1$,
\eqref{eq:allden} gives
\[
 \Bigl|y-\frac pq\Bigr|>\frac{K^{-\mu'}|Q|^{1-\mu'}}{|cx+d|\,q}\ge\frac{K^{-\mu'}C^{1-\mu'}}{|cx+d|}\,q^{-\mu'},
\]
and this exceeds $q^{-\mu}$ for large $q$.
\end{proof}

%% =====================================================================
\section{Decay on a neighbourhood}\label{sec:decay}

This section proves condition (b) of Theorem~\ref{thm:criterion} for both families, uniformly in the direction. The
method is the one of the $\zeta(5)$ paper \cite{Litsey2026c}. The determinant is an average over configurations of
points. The average is bounded by an energy problem with an external field. A single positive measure $\sigma$ then
gives an upper bound by weak duality. Two features are new. The functional is an integral over a complex contour
rather than a real interval, and the perturbation of $X$ adds point masses at the negative integers. We first prove a
general bound (Theorem~\ref{thm:dual}) and then check its hypotheses for an explicit contour and measure
(Theorem~\ref{thm:decay}).

\subsection{The field}
Put $(K,B,e,r)=(6,2,3,2)$. These are the scaled numbers of poles and zeros of $R_n$ and their orders. For $s\in\C$ put
\begin{multline*}
 \Phi(s)=r(K\log K-K)-e(B\log B-B)\\
 +e\int_0^B\log|s-x|\,dx-r\int_0^K\log|s+x|\,dx,
\end{multline*}
\[
 V(s)=\Phi(s)-2\pi|\operatorname{Im}s|,\qquad W_\sigma(s)=V(s)-2U^\sigma(s).
\]
By Stirling's formula $\log|R_n(ns)|=n\Phi(s)+O(\log n)$ away from the poles and zeros. The term $-2\pi|\operatorname{Im}s|$
is the decay of $\pi\cot\pi s$ and of its derivatives off the real axis. For $x\in[0,K]$ put
\begin{multline*}
 \omega(x)=r\bigl(K\log K-x\log x-(K-x)\log(K-x)\bigr)\\
 +e\bigl((B+x)\log(B+x)-x\log x-B\log B\bigr).
\end{multline*}
Then $\log c_k\le n\,\omega(k/n)$ for $0\le k\le6n$, by the entropy bound for binomial coefficients. Finally put
$\kappa=27/2-9\log3$ and $c(a,b)=3|a|+|b|$.

A \emph{contour} is a polyline $\Gamma$ in the half-plane $\operatorname{Re}s>0$, symmetric under conjugation. It runs from
$-i\infty$ to $+i\infty$, its two unbounded edges are vertical rays, and it meets $\R$ at one point $x_0\in(0,B)$.

\begin{theorem}\label{thm:dual}
Let $\Gamma$ be a contour. Let $\sigma$ be a positive measure on $\C$ of mass $3$ with compact support, invariant
under conjugation, with
$\sigma(D(s,r))\le A\sqrt r$ for every disc of radius $r\le r_0$, for some $A,r_0>0$. Let $W^*$, $a^*$ and $\rho$ be
real numbers with
\begin{gather*}
 W_\sigma(s)\le W^*\quad(s\in\Gamma),\qquad a^*\le W^*+\rho,\\
 \omega(x)-2U^\sigma(-x)\le a^*\quad(0\le x\le6).
\end{gather*}
Put $D=-\kappa-I(\sigma)-3W^*$. Then for every $\varepsilon>0$ there is $n_0$ such that
\[
 |J_n(a,b,a\zeta(3)+b\zeta(4)+t)|\le c(a,b)^{3n}e^{-(D-\varepsilon)n^2}
\]
for all $n\ge n_0$, all real $a,b$ with $c(a,b)>0$, and all real $t$ with $|t|\le c(a,b)e^{-\rho n}$.
\end{theorem}

The bound is uniform in $(a,b)$ and $t$. This is what the ratio family needs.

\begin{proof}
Fix $n$ and write $c=c(a,b)$. We take $X=a\zeta(3)+b\zeta(4)+t$, and we write $L_X$ for the functional at $X$.

\emph{Step 1: the contour representation.} Put $\cot_\pi(z)=\pi\cot\pi z$ and $K_{a,b}(z)=3a\cot_\pi'(z)+b\cot_\pi''(z)$. Let
$F=R_nP$ with $\deg P<6n$. Since $F=O(z^{-2})$, the sums in \eqref{eq:Lseries} converge. The residue of
$F^{(j)}\cot_\pi$ at an integer $m$ is $F^{(j)}(m)$, and on a closed contour
$\oint F^{(j)}\cot_\pi=(-1)^j\oint F\cot_\pi^{(j)}$. Let $\Gamma_n$ be the contour $n\Gamma$, moved by at most
$1$ near the crossing so that it meets $\R$ at a point $v\notin\Z$ with $0<v<2n+1$. The integers $m$ with $1\le m<v$
are zeros of $R_n$ of order $3$, so their residues vanish. Integrating over a large rectangle and letting it grow,
\[
 L_{X_0}(P)=\frac1{2\pi i}\int_{\Gamma_n}R_n(z)P(z)K_{a,b}(z)\,dz,\qquad X_0=a\zeta(3)+b\zeta(4),
\]
with $\Gamma_n$ oriented downward. By \eqref{eq:finite} the dependence on $X$ is $L_X(P)=L_{X_0}(P)+6t\sum_kc_kP(-k)$. So
$L_X$ is integration against the complex measure $\nu=\nu_\Gamma+\nu_{\rm at}$, where $\nu_\Gamma$ is the weight on
$\Gamma_n$ and $\nu_{\rm at}=6t\sum_kc_k\delta_{-k}$.

\emph{Step 2: Heine's formula.} For a complex measure $\nu$ with finite moments, Andr\'eief's identity
\cite{Forrester2018Andreief} gives
\[
 \det\Bigl[\int z^{i+j}d\nu\Bigr]_{i,j<M}=\frac1{M!}\int\!\cdots\!\int\prod_{i<j}(z_i-z_j)^2\,d\nu(z_1)\cdots d\nu(z_M).
\]
With $M=3n$ we obtain $|J_n|\le\bigl(M!\prod_{i<M}i!^2\bigr)^{-1}\int\prod_{i<j}|z_i-z_j|^2\,d|\nu|^{M}$.

\emph{Step 3: the weights.} On $\Gamma_n$ we have $|K_{a,b}(z)|\le c\,C(1+|z|)^Ce^{-2\pi|\operatorname{Im}z|}$, uniformly in
$(a,b)$. Together with Stirling's formula this gives, for $z=ns$ with $s\in\Gamma$,
\[
 d|\nu_\Gamma|(z)\le c\,Cn^C(1+|s|)^Ce^{nV(s)}\,|dz|.
\]
The atoms satisfy $6|t|c_k\le6c\,e^{-\rho n+n\omega(k/n)}$. In the scaled variable $s=z/n$, a point at an atom $-k$ has
weight at most $c\,e^{n(\omega(k/n)-\rho)}$. Two points at the same atom contribute nothing, because the Vandermonde
factor vanishes.

\emph{Step 4: the energy bound.} Let $\upsilon$ be the measure $n^{-1}\sum_i\delta_{z_i/n}$ of mass $3$. Every
configuration with $k$ points on the contour and $3n-k$ points at the atoms has integrand at most
\[
 c^{3n}n^{9n^2}\exp\Bigl(n^2\Bigl[\int\widetilde W\,d\upsilon+\iint_{s\ne u}\log|s-u|\,d\upsilon(s)\,d\upsilon(u)\Bigr]+O(n\log n)\Bigr),
\]
where $\widetilde W=V$ on $\Gamma$ and $\widetilde W(-x)=\omega(x)-\rho$ at the atoms. Now
$\int\widetilde W\,d\upsilon-I(\upsilon)=\int(\widetilde W-2U^\sigma)\,d\upsilon-I(\upsilon-\sigma)+I(\sigma)$. By the two
hypotheses and $a^*-\rho\le W^*$, we have $\widetilde W-2U^\sigma\le W^*$ on the support of $\upsilon$. The measure
$\upsilon-\sigma$ has total mass $0$. After replacing each point mass by the uniform measure on a circle of radius
$n^{-2}$, its energy is nonnegative \cite[Lemma~I.1.8]{SaffTotik}. Let $\tilde\upsilon$ be the smoothed measure and
$\delta=n^{-2}$. The mean of $\log|z-u|$ over the circle of radius $\delta$ about $u$ is $\log\max(|z-u|,\delta)$.
Hence the energy of $\tilde\upsilon$ is at most the energy of $\upsilon$ off the diagonal plus the self-energy of the
circles, which is $3n\cdot n^{-2}\log(1/\delta)=O(n^{-1}\log n)$. For the potential,
$0\le U^\sigma(s)-\int U^\sigma d(\text{circle})=\int_{|u-s|<\delta}\log(\delta/|u-s|)\,d\sigma(u)$. By
$\sigma(D(s,r))\le A\sqrt r$ this is at most $\int_0^\infty\sigma(D(s,\delta e^{-\tau}))\,d\tau\le2A\sqrt\delta=2A/n$. So
$\int U^\sigma d\tilde\upsilon\ge\int U^\sigma d\upsilon-6A/n$, and the bracket is at most $3W^*+I(\sigma)+o(1)$.

\emph{Step 5: the normalization and the tail.} The integral over the unbounded edges converges, because $W_\sigma$
tends to $-\infty$ linearly in $|\operatorname{Im}s|$ there and the factors $(1+|s|)^C$ are polynomial. The number of
choices of atoms is at most $(6n+1)^{3n}=e^{O(n\log n)}$. Finally
$2\sum_{i<3n}\log i!=9n^2\log(3n)-\tfrac{27}2n^2+O(n\log n)$. So
\begin{align*}
 \log|J_n|&\le3n\log c+n^2\bigl(3W^*+I(\sigma)+9\log n-9\log(3n)+\tfrac{27}2\bigr)+o(n^2)\\
 &=3n\log c-(D-o(1))n^2.
\end{align*}
All the error terms are uniform in $(a,b)$ and $t$.
\end{proof}

\begin{remark}
In the formal proof the contour hypothesis is stated for each large $n$ separately: the contour $\Gamma_n$ may depend on
$n$, with a bounded number of vertices, and $W_\sigma\le W^*+\eta$ is required on its pieces for every $\eta>0$. On the
two vertical rays it also requires a linear decrease of $W_\sigma$ in $|\operatorname{Im}s|$, which for a measure of
compact support is the fact used in Step~5. The statement above, with one fixed contour, is the case we use.
\end{remark}

\subsection{The contour and the measure}
We now choose $\Gamma$ and $\sigma$. The line $\operatorname{Re}s=x_0$ used by Lai gives a value of $D$ below the value
we need. A contour that bends to the right in the upper half-plane lowers the field $V$ along it, and a measure on a
segment of this contour can follow the bend.

\emph{The contour.} The upper half of $\Gamma$ is the segment from $x_0=8/5$ to $P=21/10+\frac{11}5i$, followed by the
vertical ray $x_1+iy$, $y\ge11/5$, with $x_1=21/10$. The lower half is its conjugate.

\emph{The measure.} Let $S_+=[A_0,A_1]$ with
\[
 A_0=\tfrac{161}{100}+\tfrac{11}{250}i,\qquad A_1=\tfrac{83}{40}+\tfrac{209}{100}i,
\]
a segment of the first edge of $\Gamma$. Write $M=(A_0+A_1)/2$ and $h=(A_1-A_0)/2$, and let $\omega_0$ be the arcsine
probability measure on $[-1,1]$. Put
\begin{multline*}
 H(z)=2^{-24}\bigl(107224+76785z+240184z^2+539459z^3-565327z^4\\
 +3717555z^5-9212693z^6-6236781z^7+16850383z^8\bigr),
\end{multline*}
and $f(\cos\varphi)=s_0|H(e^{i\varphi})|^2$, where $s_0>0$ is chosen so that the constant term is $c_0=3/2$. Thus
$f=\sum_{k\le8}c_kT_k$ with rational coefficients, and $f\ge0$ on $[-1,1]$. Let $s_+$ be the image of $f\,d\omega_0$
under $x\mapsto M+hx$, let $s_-$ be its conjugate, and put $\sigma=s_++s_-$. Then $\sigma$ has mass $2c_0=3$.

\emph{The constants.} $W^*=-130471/10000$, $\rho=42.94856306$.

\begin{theorem}\label{thm:decay}
For $\Gamma$, $\sigma$, $W^*$ and $\rho$ as above, the hypotheses of Theorem~\ref{thm:dual} hold, and
$I(\sigma)\le208871/10^5$. Hence $D\ge33.4401$, and for every $\varepsilon>0$ and all large $n$,
\[
 |J_n(a,b,a\zeta(3)+b\zeta(4)+t)|\le c(a,b)^{3n}e^{-(33.4401-\varepsilon)n^2}
\]
for all real $a,b$ with $c(a,b)>0$ and all $|t|\le c(a,b)e^{-42.94856306\,n}$.
\end{theorem}

The constant follows from $D=-\kappa-I(\sigma)-3W^*\ge-\tfrac{27}2+9\log3-2.08871+39.1413>33.4401$. The proof
occupies the rest of this subsection. It combines closed forms for the potentials (Lemmas~\ref{lem:V}
and~\ref{lem:joukowski}), a cover of the compact part of $\Gamma$ by $40$ cells (Proposition~\ref{prop:cells}), the energy
(Lemma~\ref{lem:energy}), the atoms (Lemma~\ref{lem:atoms}) and the tail (Lemma~\ref{lem:tail}). By symmetry under
conjugation, $W_\sigma(\bar s)=W_\sigma(s)$, so it suffices to treat the upper half of $\Gamma$. The mass bound
$\sigma(D(s,r))\le A\sqrt r$ holds because a disc of radius $r$ meets each segment in an interval of length at most $2r$,
whose arcsine mass is $O(\sqrt r)$, and $f$ is bounded.

\begin{remark}\label{rem:Dsharper}
The cell computation gives $I(\sigma)\le2.0886595$, and hence $D\ge33.44015$. The formal proof records the rounded
values $I(\sigma)\le2.08871$ and $D\ge33.4401$, which are all the proof of Theorem~\ref{thm:main} uses. The paper uses
the same values.
\end{remark}

\begin{lemma}[the field]\label{lem:V}
Put $\kappa_V=12\log3+6\log2-6$ and, with principal logarithms,
\begin{multline*}
 \widetilde\Psi(s)=e\bigl[s\log s-s-(s-B)(\log(B-s)+i\pi)+(s-B)\bigr]\\
 -r\bigl[(s+K)\log(s+K)-(s+K)-s\log s+s\bigr]+2\pi is.
\end{multline*}
\begin{enumerate}[label=(\roman*)]
\item $V(s)=\kappa_V+\operatorname{Re}\widetilde\Psi(s)$ for $\operatorname{Im}s\ge0$.
\item $\widetilde\Psi$ is analytic off $(-\infty,0]\cup[B,\infty)$, and
 $\widetilde\Psi''(s)=5/s+3/(B-s)-2/(s+K)$.
\item For $m\ge2$, $|\widetilde\Psi^{(m)}(s)|\le(m-2)!\bigl(5|s|^{1-m}+3|s-B|^{1-m}+2|s+K|^{1-m}\bigr)$, and
 $|\widetilde\Psi'(s)|\le5|\log|s||+3|\log|B-s||+2|\log|s+K||+11\pi$.
\end{enumerate}
\end{lemma}

\begin{proof}
For $\operatorname{Im}s>0$, $\int_0^B\log(s-x)\,dx=s\log s-s-(s-B)\log(s-B)+(s-B)$ with $\log(s-B)=\log(B-s)+i\pi$, and
$\int_0^K\log(s+x)\,dx=(s+K)\log(s+K)-(s+K)-s\log s+s$. Also $-2\pi\operatorname{Im}s=\operatorname{Re}(2\pi is)$. The real
parts give (i), and by continuity (i) holds on $(0,B)$. Differentiating gives
$\widetilde\Psi'=5\log s-3\log(B-s)-2\log(s+K)-\pi i$ and (ii). The bounds (iii) follow from
$|d^{j}(s-p)^{-1}/ds^j|=j!/|s-p|^{j+1}$ and $|\arg|\le\pi$.
\end{proof}

The potential of a measure $g\,d\omega_0$ on $[-1,1]$, transported to a segment, has a closed form in the Joukowski
variable. For $\zeta\notin[-1,1]$ let $\phi=\phi(\zeta)$ be the root of $\phi^2-2\zeta\phi+1=0$ with $|\phi|>1$.

\begin{lemma}[segment potentials]\label{lem:joukowski}
Let $\nu$ be the image of $g\,d\omega_0$ under $x\mapsto M_c+h_cx$, with $g=g_0+\sum_{k\ge1}g_kT_k\ge0$ a polynomial. For
$s\notin\operatorname{supp}\nu$ put $\zeta=(s-M_c)/h_c$ and
\[
 G(s)=-g_0\log h_c-g_0\log(\phi/2)+\sum_kg_k\phi^{-k}/k .
\]
\begin{enumerate}[label=(\alph*)]
\item $U^\nu=\operatorname{Re}G$, and $G$ is analytic off the segment. On the segment, at $\zeta=\cos\varphi$,
 $U^\nu=-g_0\log(|h_c|/2)+\sum_k(g_k/k)\cos k\varphi$.
\item Put $c=(|\zeta-1|+|\zeta+1|)/2$ and $d=(|\zeta+1|-|\zeta-1|)/2$. Then $|\phi|=c+\sqrt{c^2-1}$ and
 $\operatorname{Re}\phi^{-k}=|\phi|^{-k}T_k(d)$.
\item With $q=(\phi-\phi^{-1})/2$ one has $d\phi/d\zeta=\phi/q$, $dq/d\zeta=\zeta/q$ and
 $G'(s)=-\bigl(g_0+\sum_kg_k\phi^{-k}\bigr)/(h_cq)$.
\item $|G^{(m)}(s)|\le g_0(m-1)!/\operatorname{dist}(s,\operatorname{supp}\nu)^m$ for $m\ge1$, and
 $\operatorname{dist}(s,\operatorname{supp}\nu)\ge|h_c|(c-1)$.
\end{enumerate}
\end{lemma}

\begin{proof}
(a) For $x\in[-1,1]$ the classical expansion $\log|\zeta-x|=\log(|\phi|/2)-\sum_k(2/k)T_k(x)\operatorname{Re}\phi^{-k}$
holds, and $\int T_kT_j\,d\omega_0=\delta_{jk}/2$ for $k\ge1$. Integrating against $g\,d\omega_0$ and adding
$\log|h_c|$ gives $U^\nu=\operatorname{Re}G$. On the segment, $\phi=e^{i\varphi}$ by continuity. (b) Write
$\zeta=\cosh(\alpha+i\beta)$ with $\alpha>0$. Then $|\zeta\pm1|=\cosh\alpha\pm\cos\beta$, so $c=\cosh\alpha$,
$d=\cos\beta$ and $\phi=e^{\alpha+i\beta}$. (c) Differentiate $\phi=\zeta+q$ and $q^2=\zeta^2-1$. (d) Off the support,
$G^{(m)}(s)=(-1)^m(m-1)!\int(s-u)^{-m}d\nu(u)$, and $\nu$ is positive of mass $g_0$. For the distance,
$|\zeta-1|+|\zeta+1|\le2|\zeta-x|+2$ for $|x|\le1$.
\end{proof}

Lemma~\ref{lem:joukowski}(b) evaluates the potential with square roots, one logarithm and the Chebyshev recurrence.
Part (c) gives all derivatives of $G$ at a point as rational functions of $\phi$, $q$ and $h_c$. Part (d) bounds the
remainder of a Taylor expansion.

\begin{proposition}[the cells]\label{prop:cells}
$W_\sigma\le W^*$ on the segment $[x_0,P]$ and on the ray segment $\{x_1+iy:11/5\le y\le33/10\}$.
\end{proposition}

\begin{proof}
Parametrize the compact part of the upper half of $\Gamma$ by four arcs: $s=M-h\cosh w$ with
$0\le w\le146127/500000$, from $A_0$ towards the vertex; $s=M+h\cos\varphi$ with $\varphi=\pi\theta$, $0\le\theta\le1$, the
support; $s=M+h\cosh w$ with $0\le w\le459681/1000000$, from $A_1$ to $P$; and $s=x_1+iy$ with $11/5\le y\le33/10$. These
arcs cover $[x_0,P]$ and the ray segment. We cut them into $9$, $26$, $1$ and $4$ cells, with endpoints dyadic in
$\theta$, on a grid of step $2^{-12}$ in $e^{-w}$, and rational in $y$.

On each cell $t\mapsto W_\sigma(s(t))$ is real analytic, by Lemmas~\ref{lem:V} and~\ref{lem:joukowski}. We bound it by its
Taylor polynomial of order $4$ at the centre $t_c$ plus a bound for the fifth derivative. The Taylor coefficients of
$V\circ s$ come from Lemma~\ref{lem:V}(ii) and Fa\`a di Bruno's formula. Those of $U^{s_\pm}\circ s$ come from
Lemma~\ref{lem:joukowski}(c), and on the support the own potential is the trigonometric polynomial of
Lemma~\ref{lem:joukowski}(a). The fifth derivative is bounded by Lemma~\ref{lem:V}(iii) and Lemma~\ref{lem:joukowski}(d),
since the Bell polynomials have nonnegative coefficients and $|s^{(k)}|$ is bounded on each arc by $|h|$,
$|h|\cosh w$ or $1$.

All values are taken at points where every quantity is rational. On the support the centre is
$\varphi_c=2\arctan\tau$ with $\tau$ rational, so $\cos\varphi_c$ and $\sin\varphi_c$ are rational. On the other two
segment arcs the centre is $w_c=-\log\varrho$ with $\varrho$ rational, so $\cosh w_c$ and $\sinh w_c$ are rational. The
value $\phi(\zeta_c)$ is then rational in the given data, up to the square roots of Lemma~\ref{lem:joukowski}(b), which
are bounded on both sides by rationals. For the cross potential we choose a rational point $\tilde\phi$ near $\phi_c$ and
use the exact point $\tilde s=M_c+h_c(\tilde\phi+\tilde\phi^{-1})/2$. The derivatives at $s_c$ differ from those at
$\tilde s$ by at most $|s_c-\tilde s|\,k!\,g_0/d^{k+1}$, by Lemma~\ref{lem:joukowski}(d), where $d$ is a lower bound for
the distance to the support. On each of the $40$ cells the resulting upper bound for $W_\sigma$ is below $W^*$. The
logarithms and $\pi$ enter only through rational enclosures.
\end{proof}

\begin{lemma}[the tail]\label{lem:tail}
Let $R_0=2091/1000$, $y_1=33/10$ and $s=x_1+iy$ with $y\ge y_1$. Then $W_\sigma(s)\le g(y):=V(s)+6\log(y+R_0)$ and
$g'(y)\le5x_1/y-2\pi+6/(y+R_0)$. Hence $W_\sigma\le g(y_1)<W^*$ for $y\ge y_1$.
\end{lemma}

\begin{proof}
For $u$ in the support, $|s-u|\le y+|x_1-u|\le y+R_0$, so $-2U^\sigma(s)\le6\log(y+R_0)$. By Lemma~\ref{lem:V},
$\partial_yV=5\arctan(x_1/y)-3\arctan((x_1-B)/y)-2\arctan((x_1+K)/y)-2\pi$. Since $x_1>B$ the last two terms are
negative, and $\arctan z\le z$. The bound for $g'$ decreases in $y$ and is $-1.988\ldots$ at $y_1$, so $g$ decreases on
$[y_1,\infty)$. Finally $g(y_1)=-13.2158\ldots<W^*$.
\end{proof}

\begin{lemma}[the energy]\label{lem:energy}
$I(\sigma)=2I(s_+)+2I(s_+,s_-)$, with $I(s_+)=-c_0^2\log(|h|/2)+\sum_kc_k^2/(2k)$, and $I(\sigma)\le208871/10^5$.
\end{lemma}

\begin{proof}
The first identity is bilinearity and $I(s_-)=I(s_+)$. By Lemma~\ref{lem:joukowski}(a),
$I(s_+)=\int U^{s_+}ds_+$ and $\int T_kf\,d\omega_0=c_k/2$. The cross term is
$I(s_+,s_-)=\pi^{-1}\int_0^\pi g(\varphi)f(\cos\varphi)\,d\varphi$ with $g=U^{s_-}(M+h\cos\varphi)$. We integrate over the
$26$ support cells. On each, the Taylor polynomial of order $4$ of $gf$ is integrated exactly. The odd terms vanish,
and the remainder is bounded by the fifth-derivative bounds of the proof of Proposition~\ref{prop:cells}, with
$|d^j\cos k\varphi/d\varphi^j|\le k^j$.
\end{proof}

\begin{lemma}[the atoms]\label{lem:atoms}
Let $\alpha(x)=\omega(x)-2U^\sigma(-x)$ on $[0,6]$. Then $\alpha''\le-49/100$, and with $x^*=2211/500$,
\[
 \sup_{[0,6]}\alpha\le\alpha(x^*)+\frac{\alpha'(x^*)^2}{2\cdot49/100}\le W^*+42.94856306 .
\]
\end{lemma}

\begin{proof}
One computes $\omega''(x)=-r/x-r/(K-x)+e/(B+x)-e/x\le-11/24-1/3$ on $(0,6)$. Next
$-2U^\sigma(-x)=\int\log\bigl((x+\operatorname{Re}u)^2+(\operatorname{Im}u)^2\bigr)d\sigma(u)$, and the second $x$-derivative
of $\log(v^2+w^2)$ is at most $1/(4v^2)$. Here $v\ge\operatorname{Re}A_0>8/5$, so this term contributes at most
$3/(4\cdot(8/5)^2)=75/256$. Hence $\alpha''<-0.4987$. A concave function with $\alpha''\le-m$ satisfies
$\alpha(x)\le\alpha(x^*)+\alpha'(x^*)(x-x^*)-m(x-x^*)^2/2\le\alpha(x^*)+\alpha'(x^*)^2/(2m)$. The values
$\alpha(x^*)=29.90146\ldots$ and $|\alpha'(x^*)|<10^{-3}$ are computed from Lemma~\ref{lem:joukowski}(a),(b) at the
rational point $-x^*$.
\end{proof}

\begin{proof}[Proof of Theorem~\textup{\ref{thm:decay}}]
The contour satisfies the definition. The measure $\sigma$ is positive of mass $3$, conjugation-invariant, and satisfies
the mass bound. By symmetry, Proposition~\ref{prop:cells} and Lemma~\ref{lem:tail}, $W_\sigma\le W^*$ on $\Gamma$. By
Lemma~\ref{lem:atoms} the second and third hypotheses hold with $a^*=W^*+\rho$. Lemma~\ref{lem:energy} bounds
$I(\sigma)$, and Theorem~\ref{thm:dual} gives the result.
\end{proof}

\begin{corollary}\label{cor:families}
For fixed integers $a,b$ with $b\ne0$, the family $Q_n(X)=J_n(a,b,X)$ satisfies condition (b) of
Theorem~\ref{thm:criterion} at $\xi=a\zeta(3)+b\zeta(4)$ with $D=33.4401$ and $\rho=42.94856306$. The family
$f_n(y)=J_n(1,-y,0)$ satisfies it at $\xi=y^*$ with the same constants.
\end{corollary}

\begin{proof}
For $Q_n$, $c(a,b)\ge1$ is fixed, and the factor $c^{3n}$ is absorbed into $\varepsilon n^2$. For $f_n$, put
$y=y^*+t$ with $|t|\le e^{-(\rho+\varepsilon)n}$. Then $f_n(y)=J_n(1,-y,\zeta(3)-y\zeta(4)+t\zeta(4))$, by
$J_n(1,-y,0)=J_n(1,-y,X_y+t\zeta(4))$ with $X_y=\zeta(3)-y\zeta(4)$. Here $c(1,-y)=3+|y|\le6$, and
$|t\zeta(4)|\le c(1,-y)e^{-\rho n}$. Theorem~\ref{thm:decay} applies with $c\le6$.
\end{proof}

%% =====================================================================
\section{Local bounds for the denominators}\label{sec:local}

The polynomial $J_n(a,b,X)$ has rational coefficients. We bound $v_p(J_n)$ from below for every prime $p$, as a
Gauss valuation in $(a,b,X)$, and put $m_n=\prod_{p\le6n}p^{-w(n,p)}$ for the resulting lower bounds $w(n,p)$.
Then $m_nJ_n(a,b,X)$ has integer coefficients. The bounds are of three kinds. For $p>3n$ we use bases of
$\Z_{(p)}[t]_{<3n}$ adapted to the residues of the nodes (Theorems~\ref{thm:O} and~\ref{thm:B7}). For
$\sqrt{8n}<p<3n$ we prove an exact law in the first and second order (Theorems~\ref{thm:L1} and~\ref{thm:L2}). For the
remaining small primes we use Lai's bound.

\subsection{Bases and weighted determinants}
Write $B(F,G)=L_n(FG)$ for the bilinear form on $\Q[t]_{<3n}$. Its Gram matrix in the monomials has determinant
$\prod_{i<3n}i!^2\,J_n$, and the monomials form a basis of $\Z_{(p)}[t]_{<3n}$ over $\Z_{(p)}$.

\begin{lemma}\label{lem:gram}
Let $(f_x)$ be a basis of $\Z_{(p)}[t]_{<3n}$ over $\Z_{(p)}$. Then
$v_p(J_n)=v_p\det[B(f_x,f_y)]-\nu_p$, where $\nu_p=2\sum_{i<3n}v_p(i!)$. In particular this holds for
the family
\[
 e_{r,i}=(t-\rho_r)^i\prod_{r'\ne r}(t-\rho_{r'})^{\mu_{r'}}\quad(i<\mu_r),\qquad g_j=t^j\prod_{r}(t-\rho_r)^{\mu_r}\quad(j<d),
\]
whenever the $\rho_r\in\Z$ are pairwise incongruent modulo $p$ and $\sum_r\mu_r+d=3n$.
\end{lemma}

\begin{proof}
A change of basis of determinant $\delta$ multiplies the Gram determinant by $\delta^2$, and $\delta$ is a unit when
both bases span $\Z_{(p)}[t]_{<3n}$. The family above spans $\Z_{(p)}[t]_{<3n}$ by the Chinese remainder theorem in
$\F_p[t]$, since the factors $t-\rho_r$ are pairwise coprime modulo $p$. So its change of basis is unimodular.
\end{proof}

We bound the determinant through the entries. Give each basis element a \emph{potential} $\pi_x\in\frac12\Z$.

\begin{lemma}\label{lem:halfint}
Let $\cO$ be a discrete valuation ring with $v(p)=1$ and $M$ a symmetric matrix over its fraction field. Suppose the
indices split into $A$, where $\pi_x\in\Z$, and $B$, where $\pi_x\in\Z+\frac12$, with $v(M_{xy})\ge\pi_x+\pi_y$ on
$A\times A$ and $B\times B$ and $v(M_{xy})\ge\pi_x+\pi_y+\frac12$ on $A\times B$. Then
\[
 v(\det M)\ge2\sum_x\pi_x+\operatorname{corank}\overline N_{BB},
\]
where $\overline N_{BB}$ is the reduction modulo $p$ of $(p^{-\pi_x-\pi_y}M_{xy})_{x,y\in B}$. The same holds with
the roles of $A$ and $B$ exchanged.
\end{lemma}

\begin{proof}
Over $\cO[\varpi]$ with $\varpi^2=p$, put $N=D^{-1}MD^{-1}$ with $D=\operatorname{diag}(\varpi^{2\pi_x})$. The
$A\times B$ entries of $N$ lie in $\varpi\cO[\varpi]$. Multiplying the $B$ rows by $\varpi$ and the $B$ columns by
$\varpi^{-1}$ gives a matrix $Z$ over $\cO$ with $\det Z=\det N$, whose reduction is block triangular with diagonal
blocks $\overline N_{AA}$ and $\overline N_{BB}$. So the reduction has rank at most $|A|+\operatorname{rank}\overline N_{BB}$,
and the Smith form gives $v(\det Z)\ge\operatorname{corank}\overline N_{BB}$.
\end{proof}

We also use the elementary fact that $\sum_{s\in\F_p}H(s)=0$ for every $H\in\F_p[s]$ of degree at most $p-2$.

\subsection{The nodes}
In this section and the next, $p$ is the prime and $\ell$ denotes the second base-$p$ digit of a node. For a node
$k=s+\ell p$ with $0\le s<p$, write $L_n(F)=\sum_kN_k(F)$ with $N_k(F)=c_k\bigl(U_kF'(-k)+V_kF(-k)\bigr)$.
Write $6n=q_6p+r_6$ and $2n=q_2p+r_2$ with $0\le r_6,r_2<p$. The residues $s\in[0,p)$ fall into four classes, each an
interval:
\begin{center}
\renewcommand{\arraystretch}{1.15}
\begin{tabular}{@{}cccc@{}}
\toprule
class & residues & $\sigma(s)$ & $L_s$\\
\midrule
$A$ & $s\le r_6$, $s<p-r_2$ & $0$ & $q_6$\\
$B$ & $s\le r_6$, $s\ge p-r_2$ & $3$ & $q_6$\\
$C$ & $s>r_6$, $s<p-r_2$ & $2$ & $q_6-1$\\
$D$ & $s>r_6$, $s\ge p-r_2$ & $5$ & $q_6-1$\\
\bottomrule
\end{tabular}
\end{center}
The nodes of residue $s$ are $k=s+\ell p$ with $0\le\ell\le L_s$.

\begin{lemma}[node valuations]\label{lem:N}
Let $8n<p^2$ and $k=s+\ell p\le6n$.
\begin{enumerate}[label=(\roman*)]
\item $v_p(c_k)=\sigma(s)$.
\item If $\ell\ge1$, then $U_k=p^{-3}u_\ell b+O(p^{-2})$ and $V_k=p^{-4}v_{s,\ell}b+O(p^{-3})$, with $u_\ell=-2H_\ell(3)$,
 $v_{s,\ell}=-2(H_\ell(3)\delta_{s,\ell}+3H_\ell(4))$ and $\delta_{s,\ell}=5H_\ell-3H_{\ell'}-2H_{\ell''}$, where
 $\ell'=\lfloor(2n+k)/p\rfloor$ and $\ell''=\lfloor(6n-k)/p\rfloor$. The terms in $a$ are smaller by one power of $p$,
 and the term in $X$ is $p$-integral.
\item If $\ell=0$, then $v(U_k)\ge0$ and $v(V_k)\ge-1$.
\end{enumerate}
Consequently, if $F=(t+s)^mR$ with $R\in\Z_{(p)}[a,b,X][t]$, then $v(N_k(F))\ge\sigma(s)-4+m$ for $\ell\ge1$, and
$v(N_k(F))\ge\sigma(s)-1+\min(m,1)$ for $\ell=0$, with $N_k(F)=0$ if $\ell=0$ and $m\ge2$.
\end{lemma}

\begin{proof}
Since $6n<p^2$, all numbers involved have at most two digits in base $p$. By Kummer's theorem $v_p\binom{6n}k$ is the
number of carries in $k+(6n-k)$, which is $[s>r_6]$, and $v_p\binom{2n+k}k=[s\ge p-r_2]$. This gives (i). The
multiples of $p$ up to $k$, $2n+k$ and $6n-k$ are $ip$ with $i\le\ell$, $\ell'$, $\ell''$, so
$H_k(j)=p^{-j}H_\ell(j)+O(1)$ and $d_k=p^{-1}\delta_{s,\ell}+O(1)$. This gives (ii) and (iii). Finally
$F(-k)=(-\ell p)^mR(-k)$ and $F'(-k)=O(p^{m-1})$.
\end{proof}

\subsection{Lai's bound and the large primes}
The following bounds are Lai's \cite[Lemma~4.1]{Lai2026}, read for the Gauss valuation.

\begin{proposition}\label{prop:crude}
For every prime $p$, $v_p(J_n)\ge-12n\lfloor\log_p(8n)\rfloor$. If $p>6n$, then $v_p(J_n)\ge0$.
\end{proposition}

\begin{proof}
Take the binomial basis $\binom ti$, $i<3n$. Its change of basis from the monomials has determinant
$\prod_{i<3n}i!^{-1}$, so $J_n=\det\bigl[B(\binom ti,\binom tj)\bigr]$. The values of $\binom ti$ and its derivative at
the integers $-k$ have $v_p\ge-\lfloor\log_pi\rfloor$. The numbers $H_k(j)$ with $k\le8n$ have $v_p\ge-j\lfloor\log_p(8n)\rfloor$.
Each entry therefore has $v_p\ge-4\lfloor\log_p(8n)\rfloor$, and the determinant has $3n$ rows. For $p>6n$ all
quantities are $p$-integral.
\end{proof}

\begin{theorem}\label{thm:O}
Let $3n<p\le7n/2$. Then $v_p(J_n)\ge-6(p-2n)$.
\end{theorem}

\begin{proof}
Here $q_6=1$, $r_6=6n-p$, $q_2=0$ and $r_2=2n$. Put $\alpha=p-2n$. The residues $s<\alpha$ form the class $A$, those in
$[\alpha,6n-p]$ the class $B$, and the rest the class $D$. Take $\rho_r=-r$ and $\mu_r=2$ for $r<\alpha$, the other
$\mu_r=0$, and $d=7n-2p$. The potentials are $-2,-1$ for $e_{r,0},e_{r,1}$ and $0$ for the $g_j$. Then
$2\sum\pi=-6\alpha$ and $\nu_p=0$. By Lemma~\ref{lem:N} every entry has $v\ge\pi_x+\pi_y$, except the entries
$B(g_i,g_j)$ at the nodes $s+p$ with $s\in B$, which have $v\ge-1$. Their coefficient of $p^{-1}$ is a constant times
\[
 \sum_{s\in B}\hat c_s\,\Pi(-s)^2(-s)^{i+j},\qquad \Pi=\prod_{r<\alpha}(t+r)^2,
\]
where $\hat c_s$ is the unit part of $c_{s+p}$. By Anton's congruence for binomial coefficients modulo $p$
\cite[Theorem~1]{Granville97} and Wilson's reflection $m!(p-1-m)!\equiv(-1)^{m+1}$, $\hat c_s\Pi(-s)^2$ is a constant multiple of $\prod_{i<\alpha}(s-i)\prod_{j=1}^{2p-6n-1}(s+j)^2$ on $B$. This
polynomial vanishes on $A$ and $D$ and its product with $s^{i+j}$ has degree at most $p-4$. So the sum over $B$ equals
the sum over $\F_p$, which is $0$. Lemmas~\ref{lem:gram} and~\ref{lem:halfint} give the bound.
\end{proof}

Lai's bound in this range is $5n-4p$, so Theorem~\ref{thm:O} gains $7n-2p$.

\begin{theorem}\label{thm:B7}
Let $n\ge10$, $7n<2p$ and $p\le6n$. Then $v_p(J_n)\ge w_7(n,p)$, where
\[
 w_7(n,p)=\begin{cases}-2n-2p,&p\le4n,\\-18n+2p-2,&4n<p,\ 2p\le9n,\\-54n+10p-10,&2p>9n.\end{cases}
\]
\end{theorem}

These are Lai's bounds \cite[Lemma~5.4]{Lai2026}. We give a short proof by a basis with shifted roots, which works
for the polynomial in $(a,b,X)$.

\begin{proof}
Put $\alpha=p-2n$ and $A_6=6n-p$. Every node is $k=s+\ell p$ with $\ell\in\{0,1\}$, and $\ell=1$ only for $s\le A_6$.
Choose integers $D_0,t_0,N_0\ge0$ with $D_0=\min(\alpha,A_6+1)$, $t_0\le D_0$, $t_0+D_0+N_0=3n$, $t_0=D_0$ if $N_0>0$,
and $D_0+N_0\le\alpha$. Put $\rho_r=-(r+p)$ for $r<D_0$ and $\rho_r=-r$ otherwise, and let $\mu_r=2$ for $r<t_0$,
$\mu_r=1$ for $t_0\le r<D_0+N_0$, and $\mu_r=0$ otherwise. Give $e_{r,i}$ the potential $i-2$ for $r<D_0$, and
potential $1$ for $D_0\le r<D_0+N_0$.

The root $-(r+p)$ sits on the node $k=r+p$ itself. If $F$ has a zero of order at least $2$ there, then
$F(-k)=F'(-k)=0$ and $N_k(F)=0$. This replaces Lai's rank condition at these nodes. At the node $k=r$ the factor
$(t+r+p)^m$ has value $p^m$, so $v(N_r(F))\ge\sigma(r)+m-1$. For $A_6<s<\alpha$ all harmonic numbers in $V_s$ have
indices below $p$, so $v(V_s)\ge0$ and $v(N_s(F))\ge\sigma(s)=2$. With Lemma~\ref{lem:N} for the other nodes, every
entry satisfies $v\ge\pi_x+\pi_y$. Lemmas~\ref{lem:gram} and~\ref{lem:halfint} give
$v_p(J_n)\ge2\sum\pi=2(N_0-t_0-2D_0)$, since $\nu_p=0$. The choices
\begin{center}
\renewcommand{\arraystretch}{1.15}
\begin{tabular}{@{}lccc@{}}
\toprule
range & $D_0$ & $t_0$ & $N_0$\\
\midrule
$p\le4n$ & $\alpha$ & $5n-p$ & $0$\\
$4n<p$, $2p\le9n+1$ & $A_6+1$ & $p-3n-1$ & $0$\\
$2p\ge9n+2$ & $A_6+1$ & $D_0$ & $2p-9n-2$\\
\bottomrule
\end{tabular}
\end{center}
give $2(N_0-t_0-2D_0)=w_7(n,p)$ in the three ranges. For $2p=9n+1$ the bound obtained is larger than $w_7$.
\end{proof}

\subsection{The law in the range \texorpdfstring{$\sqrt{8n}<p<3n$}{sqrt(8n) < p < 3n}}
Lai's bound in this range \cite[Lemma~4.4]{Lai2026} is a greedy count. The $j$-th index of the residue $s$ costs the
\emph{increment} $2j-6+\sigma(s)$, plus $3$ if $j>2L_s$, and Lai's bound is the sum of the $3n$ smallest increments
minus $\nu_p$. We call this sum $\Lex(n,p)$. Let $\lambda$ be the largest increment taken and $d$ the number of
increments taken that equal $\lambda$. Let $\mu_s$ be the number of indices of $s$ with increment below $\lambda$, so
that $d=3n-\sum_s\mu_s$. The two classes with $\sigma\equiv\lambda\pmod2$ are the \emph{tie classes}: $(A,C)$ if $\lambda$
is even and $(B,D)$ if $\lambda$ is odd. Let $Z$ be the smaller tie class, $W$ the larger, and $\mathcal N$ the union
of the other two classes. Put
\[
 x_1=\max\bigl(0,\min(d-|Z|,|W|-d)\bigr),\qquad x_2=\max(0,x_1-|\mathcal N|).
\]

\begin{lemma}[shape of the greedy]\label{lem:S}
Let $\sqrt{8n}<p<3n$. Then $\lambda\ge0$, every residue $r$ satisfies $2\mu_r+\sigma(r)=\lambda+4+[r\in\mathcal N]$, and
every tie residue satisfies $\mu_r+1\le2L_r$. Moreover
\begin{equation}\label{eq:lemT}
 2d+|\mathcal N|+2=p(2q_6-3q_2-2-\lambda).
\end{equation}
If $x_2>0$, then $\lambda=2q_6-3q_2-3$, $2d+2=|Z|+|W|$, $x_1=d-|Z|$ and $x_2=x_1-|\mathcal N|$.
\end{lemma}

\begin{proof}
The indices with negative increment are $j\le2$ on $A$ and $j=1$ on $B$ and $C$. There are $r_6+1+p-r_2<3n$ of
them, as one checks in the three ranges $p\le3n/2$, $3n/2<p\le2n$ and $2n<p<3n$. So $\lambda\ge0$. The increment of the
index $2L_c$ of a class $c$ is $T_c=4L_c-6+\sigma_c$, and the smallest is $T_C=4q_6-8$. At least
$12n-2p-r_6-r_2+1\ge3n$ indices have increment at most $4q_6-7$, so $\lambda\le T_c+1$ for every class. On a tie class
the parity then gives $\lambda\le T_c$. Hence the cap $j>2L$ is never reached, and the classes are filled up to
$\lambda$, which is the formula for $\mu_r$. Summing it over $r$, with
$\sum_r\sigma(r)=2(p-1-r_6)+3r_2$ and $3r_2-r_6=p(q_6-3q_2)$, gives \eqref{eq:lemT}. Write
$k=2q_6-3q_2-2-\lambda\ge1$. If $x_2>0$ and $k\ge2$, then $2|Z|+3|\mathcal N|<2$, so all residues lie in one class.
That class contains $0$, hence is $A$, and then $r_6=p-1$ and $r_2=0$, which contradicts $r_6\equiv3r_2\pmod p$. So $k=1$
and the remaining identities follow.
\end{proof}

The unit parts of the weights are controlled by a product formula. By Anton's congruence
\cite[Theorem~1]{Granville97} the unit part of $c_{s+\ell p}$ is $\kappa_c(\ell)\hat c_s$ modulo $p$, where $\kappa_c(\ell)$
depends only on the class $c$ and on $\ell$, and $\hat c_s$ is a product of factorials of the low digits.

\begin{lemma}\label{lem:U}
Put $w_s=\hat c_s\prod_{r\ne s}(r-s)^{2\mu_r}$ and $Q(s)=\prod_{\beta\in\mathcal N}(s-\beta)$. On each tie class $c$
there is $K_c\in\F_p^\times$ with $w_s\equiv K_cQ(s)$.
\end{lemma}

\begin{proof}
Inside a class, $\hat c_{s+1}/\hat c_s\equiv(s-r_6)^2(s-b_\varepsilon)^3(s-p+1)^{-5}$ with $b_\varepsilon=p-r_2-1$. The ratio of
the products is $\prod_{r\ne s}(r-s)^{2(\mu_{r+1}-\mu_r)}$, and the exponent is nonzero only at the class boundaries
$r_6$, $b_\varepsilon$ and $p-1$. By Lemma~\ref{lem:S}, $2(\mu_{b+1}-\mu_b)=-(\sigma_{b+1}-\sigma_b)+([b+1\in\mathcal N]-[b\in\mathcal N])$.
The jumps of $\sigma$ cancel the exponents from $\hat c$, and what remains is
$\prod_b(s-b)^{[b+1\in\mathcal N]-[b\in\mathcal N]}=Q(s+1)/Q(s)$.
\end{proof}

\begin{theorem}\label{thm:L1}
For every prime $\sqrt{8n}<p<3n$, $v_p(J_n)\ge\Lex(n,p)+x_1$.
\end{theorem}

\begin{proof}
Take the family of Lemma~\ref{lem:gram} with $\rho_r=-r$, the multiplicities $\mu_r$ and $d$ as above. Give $e_{r,i}$
the potential $i-2+\sigma(r)/2$ and each $g_j$ the potential $\lambda/2$. Then $2\sum\pi$ is the sum of the $3n$ smallest
increments, by Lemma~\ref{lem:S}.

\emph{Entries.} By Lemmas~\ref{lem:N} and~\ref{lem:S}, every entry has $v\ge\pi_x+\pi_y$, and equality can occur only
for two elements of one residue, or for a pair involving the $g_j$ at a tie residue. The elements of tie classes and
the $g_j$ have potentials of the same parity, and Lemma~\ref{lem:halfint} applies with this group in the role of $B$.

\emph{The reduced form.} By Lemma~\ref{lem:N}, the reduction of the tie part is
$\overline N_T=\sum_{s\in Z\cup W}w_sE_s^{\top}H_{c(s)}E_s$. Here $E_s$ maps a coefficient vector $(x,h)$ to
$(x_{s,0},\dots,x_{s,\mu_s-1},h(s))$, with $h(s)=\sum_jh_j(-s)^j$, and $H_c$ is a Hankel matrix
$(\Phi_c(i+i'))_{i,i'\le\mu_c}$ of class constants. We replace the numbers $\Phi_c(m)$ of each tie class by independent
indeterminates. The actual matrix is a specialization of this generic matrix, and the rank does not increase under
specialization. So it suffices to find $x_1$ independent kernel vectors of the generic matrix.

\emph{Kernel vectors.} Let $h\in\F[s]$ vanish on $Z$, with $\deg h\le\min(d-1,|Z|+|W|-d-1)$. For $s\in W$ put
$x_s=-H_\mu^{-1}\varphi\,h(s)$, where $H_\mu$ is the leading $\mu\times\mu$ block of $H_W$ and $\varphi=(\Phi_W(i+\mu))_{i<\mu}$,
and put $x=0$ elsewhere. Then $H_WE_s(x,h)=(0,\kappa_Wh(s))$ with $\kappa_W=\Phi_W(2\mu)-\varphi^{\top}H_\mu^{-1}\varphi$. The
pairing of $(x,h)$ with any $(x',h')$ is
\[
 \kappa_W\sum_{s\in W}w_sh(s)h'(s)=\kappa_WK_W\sum_{s\in\F_p}Q(s)h(s)h'(s),
\]
by Lemma~\ref{lem:U}, because $h=0$ on $Z$ and $Q=0$ on $\mathcal N$. The polynomial $Qhh'$ has degree at most
$|\mathcal N|+(|Z|+|W|-d-1)+(d-1)=p-2$, so the sum vanishes. The admissible $h$ form a space of dimension $x_1$.

Lemmas~\ref{lem:gram} and~\ref{lem:halfint} now give $v_p(J_n)\ge2\sum\pi+x_1-\nu_p$.
\end{proof}

%% =====================================================================
\section{The second level of the law}\label{sec:second}

\begin{theorem}\label{thm:L2}
For every prime $\sqrt{8n}<p<3n$, $v_p(J_n)\ge\Lex(n,p)+x_1+x_2$.
\end{theorem}

If $x_2=0$ this is Theorem~\ref{thm:L1}, so we assume $x_2>0$. By Lemma~\ref{lem:S},
\begin{equation}\label{eq:deg}
 \lambda=2q_6-3q_2-3,\qquad 3|\mathcal N|+2|Z|+2x_2-2=p-4 .
\end{equation}
The idea is as follows. In Theorem~\ref{thm:L1} the gain came from a kernel of the reduced form. Here we need one more
unit of valuation for $x_2$ elements. We raise their potentials, and we show that the entries still meet the
required bounds. Most pairs meet them by the node valuations, if the basis elements at a residue of $W$ are built
from local orthogonal polynomials. The remaining pairs need a sum over $W$ to vanish modulo $p$ or modulo $p^2$. The
second of these sums is a Faulhaber sum plus a sum of harmonic numbers. The two have opposite totals by
\eqref{eq:deg}, and they cancel.

Let $W$ have the digit flags $\beta_W=[W\subseteq C\cup D]$ and $\varepsilon_W=[W\subseteq B\cup D]$, and write
$\sigma=\sigma_W$, $L=L_W$, $\mu=\mu_W$ and $w=|W|$. Put $\mu'_r=\mu_r+1$ for $r\notin W$ and $\mu'_r=\mu$ on $W$.

\subsection{Generic node constants}
By Lemma~\ref{lem:N}, the nodes $k=s+\ell p$ with $s\in W$ and $1\le\ell\le L$ have leading constants
$\tau_\ell=(\kappa_\ell u_\ell,\kappa_\ell v_\ell)$, independent of $s\in W$. We replace them by indeterminates. Let
$T=(T^u_\ell,T^v_\ell)_{1\le\ell\le L}$ and define $L_T$ by adding
$b\,p^{\sigma-4}\mathrm{cl}_k\bigl[(T^u_\ell-\kappa_\ell u_\ell)\,pF'(-k)+(T^v_\ell-\kappa_\ell v_\ell)F(-k)\bigr]$ to $N_k(F)$
at these nodes, where $\mathrm{cl}_k$ is a fixed integer lift of the unit part of $c_k/\kappa_\ell$. Let $J_n(T)$ be the
determinant for $L_T$. At $T=\tau$ it is $J_n$, and $\tau$ is $p$-integral, so $v_p(J_n)\ge v_p(J_n(T))$ for the Gauss
valuation in $(a,b,X,T)$. From now on $\cO$ is the valuation ring of this Gauss valuation on $\Q_p(a,b,X,T)$.

\begin{lemma}\label{lem:H}
Let $1\le L<p$ and $0\le\mu\le2L$, and let $\cL_T(\Phi)=\sum_{\ell=1}^L\bigl(T^v_\ell\Phi(-\ell)+T^u_\ell\Phi'(-\ell)\bigr)$.
Then $\det(\cL_T(y^{i+j}))_{i,j<\mu}$ is a nonzero element of $\F_p[T]$.
\end{lemma}

\begin{proof}
Specialize. If $\mu\le L$, put $T^v_\ell=1$ for $\ell\le\mu$ and all other $T=0$. If $\mu>L$, put $T^u_\ell=1$ for
$\ell\le\mu-L$, $T^v_\ell=1$ for $\mu-L<\ell\le L$, and all other $T=0$. In both cases the form
$(\phi,\psi)\mapsto\cL(\phi\psi)$ on $\F_p[y]_{<\mu}$ is, by the Chinese remainder theorem, an orthogonal sum of the
nondegenerate local forms $\phi\psi(-\ell)$ and $(\phi\psi)'(-\ell)$, since the nodes $-\ell$ are distinct modulo $p$.
\end{proof}

So there is a polynomial $A(y)=\sum_{j\le\mu}A_jy^j$ over $\cO$, with leading coefficient $A_\mu=\Delta$ a unit of $\cO$
(the Hankel determinant of Lemma~\ref{lem:H}), such that $\cL_T(y^iA)=0$ for $i<\mu$. For $s\in W$ put
\[
 P_s(t)=\sum_{j\le\mu}A_jp^{\mu-j}(t+s)^j .
\]
Then $P_s(-s-\ell p)=p^\mu A(-\ell)$ and $pP_s'(-s-\ell p)=p^\mu A'(-\ell)$. The polynomial $P_s$ is a root-shifted
form of $\Delta(t+s)^\mu$: it is congruent to it modulo $p$.

\subsection{The basis}
Put $\Pi=\prod_{s\in W}P_s\prod_{r\notin W}(t+r)^{\mu'_r}$ and take
\[
 e_{r,i}=\frac{(t+r)^i\Pi}{(t+r)^{\mu'_r}}\ (r\notin W,\ i<\mu'_r),\qquad e_{s,i}=\frac{(t+s)^i\Pi}{P_s}\ (s\in W,\ i<\mu),
\]
and $g_j=t^j\Pi$ for $j<x_2$.
The count is $\sum_r\mu_r+|Z|+|\mathcal N|+x_2=3n$, by Lemma~\ref{lem:S}. Since $P_s\equiv\Delta(t+s)^\mu$ modulo $p$,
this is a basis of $\cO[t]_{<3n}$ with unimodular change of basis, as in Lemma~\ref{lem:gram}. The potentials are
$\pi(e_{r,i})=i-2+\sigma(r)/2$ and $\pi(g_j)=\lambda/2+1$. By Lemma~\ref{lem:S} and \eqref{eq:deg},
\begin{align*}
 2\sum\pi&=\sum_{\text{$3n$ smallest}}(\text{increments})+|Z|\lambda+|\mathcal N|(\lambda+1)+x_2(\lambda+2)-d\lambda\\
 &=\Lex(n,p)+\nu_p+x_1+x_2,
\end{align*}
using $d=|Z|+|\mathcal N|+x_2$. So it suffices to show $v(L_T(f_xf_y))\ge\pi_x+\pi_y$ for every pair, with the
half-integral refinement of Lemma~\ref{lem:halfint} between the two parity groups, and corank $0$.

\subsection{Node jets}
For $F\in\cO[t]$ and a node $k$ put
\[
 \nu_k(F)=\min\bigl(v F(-k),\,v(pF'(-k))\bigr).
\]
The map
$F\mapsto(F(-k),pF'(-k))$ is multiplicative for the product
\[
 (f_0,f_1)(g_0,g_1)=(f_0g_0,\,f_0g_1+f_1g_0),
\]
so
$\nu_k(FG)\ge\nu_k(F)+\nu_k(G)$. Also $\nu_k(t+r)\ge[r\equiv k]$, and $\nu_k(P_s)\ge\mu[s\equiv k]$.

\begin{lemma}[harmonic digits]\label{lem:Hdig}
Let $m=cp+e$ with $0\le e<p$ and $c<p$. Then $H_m-p^{-1}H_c-H_e\in p\Z_{(p)}$. For $k=\ell p+s\le6n$,
\[
 d_k+p^{-1}\delta_{s,\ell}+\lambda_0(s)\in p\Z_{(p)},\qquad
 \lambda_0(s)=3H_{r_2+s-\varepsilon p}+2H_{r_6-s+\beta p}-5H_s,
\]
with $(\beta,\varepsilon)=([s>r_6],[s\ge p-r_2])$.
\end{lemma}

\begin{proof}
The terms $i=jp$ of $H_m$ give $p^{-1}H_c$. Every other term satisfies $1/(jp+i)-1/i\in p\Z_{(p)}$, which leaves
$cH_{p-1}+H_e$, and $H_{p-1}\in p\Z_{(p)}$. Apply this to $k$, $2n+k$ and $6n-k$, whose digits are $(\ell,s)$,
$(q_2+\ell+\varepsilon,r_2+s-\varepsilon p)$ and $(q_6-\ell-\beta,r_6-s+\beta p)$.
\end{proof}

By the same digit expansion of the factorials, $c_k=p^\sigma\kappa_\ell\mathrm{cl}_k$ modulo $p^{\sigma+2}$ with
$\mathrm{cl}_k\equiv\bar c_s(1+p\ell\lambda_0(s))$ modulo $p^2$, where $\bar c_s$ is an integer lift of
$\hat c_s(1+pe_0(s))$ and $e_0(s)=-2(q_6-\beta)H_{r_6-s+\beta p}+3(q_2+\varepsilon)H_{r_2+s-\varepsilon p}$.

\begin{lemma}[node jets]\label{lem:J}
Let $k=\ell p+s$ with $\ell\ge1$, $s\in W$, and $F\in\cO[t]$ with $\nu_k(F)\ge e$. Put $f_0=F(-k)$ and $f_1=pF'(-k)$.
Then the node term of $L_T$ at $k$ satisfies
\begin{multline*}
 N^T_k(F)\equiv b\,p^{\sigma-4}\mathrm{cl}_k\bigl[T^u_\ell f_1+T^v_\ell f_0\bigr]-b\,p^{\sigma-3}\kappa_\ell\mathrm{cl}_k\lambda_0(s)u_\ell f_0\\
 +a\,p^{\sigma-3}\kappa_\ell\mathrm{cl}_k\bigl[u^a_\ell f_1+v^a_\ell f_0\bigr]
\end{multline*}
modulo $p^{\sigma-2+e}$, with $u^a_\ell=-3H_\ell(2)$ and $v^a_\ell=-3(\delta_{s,\ell}H_\ell(2)+2H_\ell(3))$. At the other
nodes with $\ell\ge1$, $v(N_k(F))\ge\sigma(s)-4+\nu_k(F)$. At the nodes $k=s<p$, $v(N_s(F))\ge\sigma(s)-1+\nu_s(F)$.
\end{lemma}

\begin{proof}
$N_k(F)=c_k(U_kp^{-1}f_1+V_kf_0)$. By Lemma~\ref{lem:Hdig},
$U_k=p^{-3}(bu_\ell+pau^a_\ell)+O(1)$ and $V_k=p^{-4}(bv_\ell+pav^a_\ell)-p^{-3}b\lambda_0u_\ell+O(p^{-2})$. Add the
modification of $L_T$. The $b$-terms then carry the factor $c_k/p^\sigma-\kappa_\ell\mathrm{cl}_k\equiv0$ modulo $p^2$.
For $\ell=0$, $v(U_s)\ge0$ and $v(V_s)\ge-1$.
\end{proof}

\subsection{Pairs at the residues of \texorpdfstring{$W$}{W}}
For an element $X$ of the basis and a residue $s$ let $\operatorname{ord}_s(X)$ be $i$ if $X=e_{s,i}$ and $\mu'_s$
otherwise, and put $\delta_X(s)=\pi_X-\operatorname{ord}_s(X)-(\sigma(s)-4)/2$. By Lemma~\ref{lem:J} a pair $(X,Y)$ meets
its bound at all nodes of residue $s$ if $\delta_X(s)+\delta_Y(s)\le0$ (same parity group) or $\le-\frac12$ (different
groups). By Lemma~\ref{lem:S}, $\delta_X(s)$ is $\pi_X-\lambda/2-1$ at $s\in Z$, $\pi_X-\lambda/2-\frac32$ at
$s\in\mathcal N$, and $\pi_X-\lambda/2$ at $s\in W$ when $X$ is not an element of $s$. The largest potentials are
$\lambda/2+1$ for the $g_j$, $\lambda/2+\frac12$ for the top elements of $\mathcal N$, $\lambda/2$ for the top elements of
$Z$, and $\lambda/2-\frac12$ for the second elements of $\mathcal N$. All others are at most $\lambda/2-1$. A direct
check gives:
\begin{itemize}
\item at the residues of $Z$ and $\mathcal N$ every pair meets its bound;
\item at the residues of $W$ the only pairs that do not are (a) an element of $s$ with $g_j$ or with a top element of
 $\mathcal N$, with deficit $1$; (b) $g_j$ with a top element of $Z$ or a second element of $\mathcal N$, and two top
 elements of $Z$ and $\mathcal N$ or of $\mathcal N$, with deficit $1$; (c) $g_j$ with $g_{j'}$ or with a top element of
 $\mathcal N$, with deficit $2$.
\end{itemize}
The nodes $k=s<p$ have room $3$, so only the nodes with $\ell\ge1$ matter.

\begin{lemma}[own pairs]\label{lem:own}
Let $s\in W$, $i<\mu$, and let $Y$ be a basis element with the factor $P_s$. Then
$\sum_{\ell}N^T_{s+\ell p}(e_{s,i}Y)\equiv0$ modulo $p^{\sigma-3+\mu+i}$.
\end{lemma}

\begin{proof}
Write $e_{s,i}Y=(t+s)^iP_sR$. The jet at the node $\ell$ is
\[
 p^{\mu+i}\bigl((y^iA)(-\ell),\,(y^iA)'(-\ell)\bigr)R(-s)
\]
plus higher terms. By Lemma~\ref{lem:J} the sum is $b\,p^{\sigma-4+\mu+i}\bar c_sR(-s)\cL_T(y^iA)$ modulo
$p^{\sigma-3+\mu+i}$, and $\cL_T(y^iA)=0$.
\end{proof}

For the pairs (b) and (c), write $XY=P_s^2R_s$ for $s\in W$. Then $XY=\Pi^2\phi$ for a rational function $\phi$, and we
put $\psi(x)=\phi(-x)$ and $G=\hat Q\psi$, with $\hat Q(x)=\prod_{r\notin W}(x-r)^{m_r}$, where $m_r=3$ on $\mathcal N$ and
$2$ on $Z$. Then $G$ is a polynomial, of degree at most $3|\mathcal N|+2|Z|+2x_2-2=p-4$ by \eqref{eq:deg}. It vanishes
at every residue outside $W$, to order at least $2$ for the pairs (c).

\begin{lemma}[unit parts on $W$]\label{lem:R}
For $s\in W$ put $f(s)=\hat c_s\prod_{r\ne s}(r-s)^{2\mu'_r}$. There is a unit $K$ with
$f(s)\equiv K\hat Q(s)\bigl(1+p\,\Xi(s)\bigr)$ modulo $p^2$, where
$\Xi(s)=(\lambda-1)H_s-2\beta_WH_{r_6+p-s}+3(1-\varepsilon_W)H_{r_2+s}$.
\end{lemma}

\begin{proof}
Put $\beta=\beta_W$, $\varepsilon=\varepsilon_W$ and $P_c(s)=\prod_{c\le r<p,\,r\ne s}(s-r)$. By Lemma~\ref{lem:S},
$2\mu'_r+\sigma(r)=\lambda+4+m_r$ for every residue $r$, with $m_r=0$ on $W$. Since
$\sigma(r)=2[r>r_6]+3[r\ge p-r_2]$, this gives the exact identity
\[
 f(s)\,P_{r_6+1}(s)^2P_{p-r_2}(s)^3=\hat c_s\,\hat Q(s)\,P_0(s)^{\lambda+4}\qquad(s\in W).
\]
Every factor is now a product of factorials of numbers in $[0,p)$. We use
$m!(p-1-m)!\equiv(-1)^m(p-1)!(1+pH_m)$ modulo $p^2$, which follows from
$\prod_{i=1}^{p-1-m}(p-i)\equiv(-1)^{p-1-m}(p-1-m)!(1-pH_{p-1-m})$ and $H_{p-1-m}\equiv H_m$. Write $u=(p-1)!$ and
$\hat c_s=\bigl(r_6!/(s!\,(r_6+\beta p-s)!)\bigr)^2\bigl((r_2+s-\varepsilon p)!/(s!\,r_2!)\bigr)^3$.

First, $P_0(s)=(-1)^{p-1-s}s!(p-1-s)!\equiv u(1+pH_s)$. Next, if $\beta=0$ then
$P_{r_6+1}(s)=\pm(p-1-s)!/(r_6-s)!$ with a sign independent of $s$. If $\beta=1$ then
$P_{r_6+1}(s)=\pm(s-r_6-1)!\,(p-1-s)!$, and the reflection with $m=r_6+p-s$ gives
$(r_6+p-s)!\,(s-r_6-1)!\equiv\pm u(1+pH_{r_6+p-s})$. In both cases
\begin{multline*}
 \Bigl(\frac{r_6!}{s!\,(r_6+\beta p-s)!}\Bigr)^2P_{r_6+1}(s)^{-2}\\
 \equiv\mathrm{const}\cdot\bigl(1-2\beta pH_{r_6+p-s}\bigr)\,\bigl(s!(p-1-s)!\bigr)^{-2}.
\end{multline*}
In the same way, if $\varepsilon=1$ then $P_{p-r_2}(s)=(s-p+r_2)!\,(-1)^{s}(p-1-s)!$, and if $\varepsilon=0$ the
reflection with $m=r_2+s$ applies, so that
\begin{multline*}
 \Bigl(\frac{(r_2+s-\varepsilon p)!}{s!\,r_2!}\Bigr)^3P_{p-r_2}(s)^{-3}\\
 \equiv\mathrm{const}\cdot(-1)^s\bigl(1+3(1-\varepsilon)pH_{r_2+s}\bigr)\,\bigl(s!(p-1-s)!\bigr)^{-3}.
\end{multline*}
Finally $(s!(p-1-s)!)^{-5}\equiv(-1)^su^{-5}(1-5pH_s)$, and the two signs $(-1)^s$ cancel. Multiplying,
$f(s)/\hat Q(s)\equiv K\bigl(1+p\,\Xi(s)\bigr)$ with $\Xi=(\lambda+4-5)H_s-2\beta H_{r_6+p-s}+3(1-\varepsilon)H_{r_2+s}$
and $K$ a unit independent of $s$.
\end{proof}

We need two facts about sums over a complete residue system. First, for $G\in\Z_{(p)}[x]$ of degree at most $p-4$,
\begin{equation}\label{eq:faulhaber}
 \sum_{s<p}G(s)\equiv p\int G\pmod{p^2},\qquad \int G:=\sum_eg_eB_e,
\end{equation}
where $G=\sum_eg_ex^e$ and $B_e$ are the Bernoulli numbers with $B_1=-\frac12$. This is Faulhaber's formula. Second,
$\sum_{m<p}m^eH_m\equiv B_e\pmod p$ for $e\le p-4$. Hence, if $G$ vanishes to order $2$ at every residue outside $W$ and
$x(s)$ is $s+\alpha$ or $c-s$ with $x(W)\subseteq[0,p)$, then
\begin{equation}\label{eq:shift}
 \sum_{s\in W}G(s)H_{x(s)}\equiv\int G\pmod p .
\end{equation}
For $x=c-s$ this follows from $H_{p-1-m}\equiv H_m$. For $x=s+\alpha$, the residues $s+\alpha$ run over a complete
system as $s$ runs over $W$ and its complement, the terms off $W$ vanish, and the shift
$\int G(x+1)=\int G+G'(0)$ changes $\int G$ only by values of $G'$ at residues outside $W$, which vanish.

\begin{lemma}[window sums]\label{lem:window}
Let $(X,Y)$ be a pair of type (b) or (c) at the residues of $W$. Then
$v\bigl(\sum_{s\in W}\sum_\ell N^T_{s+\ell p}(XY)\bigr)\ge\sigma+2\mu-4+\theta$, where $\theta=1$ for type (b) and
$\theta=2$ for type (c).
\end{lemma}

\begin{proof}
For $\ell\ge1$ the jets of $F=P_s^2R_s$ are $f_0=p^{2\mu}A(-\ell)^2R_s(-k)$ and
$f_1=p^{2\mu}\bigl(2AA'(-\ell)R_s(-k)+pA(-\ell)^2R_s'(-k)\bigr)$. Insert these into Lemma~\ref{lem:J}, together with
$R_s(-k)\equiv R_s(-s)-\ell pR_s'(-s)$ and $\mathrm{cl}_k\equiv\bar c_s(1+p\ell\lambda_0)$. Summing over $\ell$, with
$(yA^2)'=A^2+2yAA'$, gives modulo $p^{\sigma+2\mu-2}$
\begin{multline*}
 p^{\sigma+2\mu-4}\bar c_s\Bigl\{b\bigl[R_s(-s)\cL_T(A^2)+p(R_s'(-s)-\lambda_0R_s(-s))\cL_T(yA^2)\\
 +p\lambda_0R_s(-s)\cL^u_{T-\tau}(A^2)\bigr]+paR_s(-s)\cL^a(A^2)\Bigr\},
\end{multline*}
with three local functionals independent of $s$. By Lemma~\ref{lem:R} and $P_r(-s)\equiv\Delta(r-s)^\mu$,
\begin{multline*}
 \bar c_sR_s(-s)\equiv KG(s)\bigl\{\Delta^{2w-2}(1+p(\Xi+e_0)(s))\\
 +2p\Delta^{2w-3}A_{\mu-1}\sigma_W(s)\bigr\}\pmod{p^2},
\end{multline*}
with $\sigma_W(s)=\sum_{r\in W,r\ne s}1/(r-s)=H_{w_1-s}-H_{s-w_0}$ for $W=[w_0,w_1]$, and $\bar c_sR_s'(-s)$ is $K\Delta^{2w-2}$
times a polynomial that vanishes off $W$, modulo $p$.

For type (b) it suffices to work modulo $p^{\sigma+2\mu-3}$, where the sum is a unit times $\sum_{s\in W}G(s)$. Since
$G$ vanishes at the residues outside $W$, this is $\sum_{s\in\F_p}G(s)=0$ in $\F_p$.

For type (c) the sums $\sum_WG$, $\sum_WG\sigma_W$, $\sum_WG\lambda_0$ and the sum of the derivative terms must be
divisible by $p$, and $\sum_WG(1+p(\Xi+e_0))$ by $p^2$. The first and the derivative sum vanish as before. By
\eqref{eq:shift}, $\sum_WG\sigma_W\equiv(1-1)\int G$ and $\sum_WG\lambda_0\equiv(3+2-5)\int G$, both $0$. By
\eqref{eq:faulhaber} and \eqref{eq:shift},
\[
 \sum_{s\in W}G(s)\bigl(1+p(\Xi+e_0)(s)\bigr)\equiv p\bigl(1+\omega\bigr)\int G\pmod{p^2},
\]
where $\omega$ is the total weight of the harmonic numbers in $\Xi+e_0$:
$\omega=(\lambda-1)-2\beta_W+3(1-\varepsilon_W)-2(q_6-\beta_W)+3(q_2+\varepsilon_W)=\lambda+2-2q_6+3q_2$. By \eqref{eq:deg},
$1+\omega=0$.
\end{proof}

\begin{proof}[Proof of Theorem~\textup{\ref{thm:L2}}]
Every entry $L_T(XY)$ is a sum of node terms. By Lemma~\ref{lem:J} and the check above, every node term meets
$\pi_X+\pi_Y$, except the nodes of $W$ for the pairs (a), (b) and (c), whose sums over $W$ meet it by
Lemmas~\ref{lem:own} and~\ref{lem:window}. Lemma~\ref{lem:halfint}, applied with corank $0$, gives
$v(\det)\ge2\sum\pi$, and Lemma~\ref{lem:gram} gives $v_p(J_n(T))\ge\Lex(n,p)+x_1+x_2$.
\end{proof}

%% =====================================================================
\section{The prime sum}\label{sec:sum}

We now sum the local bounds. For $n\ge10$ and a prime $p$ put
\[
 w(n,p)=\begin{cases}0,&p>6n,\\ w_7(n,p),&7n<2p\le12n,\\ -6(p-2n),&3n<p,\ 2p\le7n,\\
 \max\bigl(\Lex(n,p)+x_1+x_2,\,-12n\lfloor\log_p8n\rfloor\bigr),&\sqrt{8n}<p<3n,\\ -12n\lfloor\log_p8n\rfloor,&\text{otherwise,}\end{cases}
\]
and $m_n=\prod_{p\le6n}p^{-w(n,p)}$. By Proposition~\ref{prop:crude} and Theorems~\ref{thm:O}, \ref{thm:B7}
and~\ref{thm:L2}, $v_p(m_nJ_n)\ge0$ for every prime $p$. So $m_nJ_n(a,b,X)$ has integer coefficients.

\subsection{The law on 22 pieces}
For $\sqrt{8n}<p<3n$ write $q=q_6=\lfloor6n/p\rfloor\ge2$ and $t=\fr{6n/p}$. The class sizes are affine in $p$ and
$n$ with bounded error, and so are $d$, $x_1$ and $x_2$. The next proposition is the result of this bookkeeping.

\begin{proposition}\label{prop:bridge}
Let $n\ge10$ and $\sqrt{8n}<p<3n$. Let $(r,[t_0,t_1])$ be the row of Table~\ref{tab:pieces} with $r=q\bmod6$ and
$t\in[t_0,t_1)$. Then
\[
 \Lex(n,p)+x_1+x_2\ge p\,c_0(q)+n\,c_1(q)-1000(q^2+1).
\]
\end{proposition}

\begin{table}[ht]
\centering
\renewcommand{\arraystretch}{1.12}
\begin{tabular}{@{}ccll@{}}
\toprule
$r$ & $t$ & $c_0(q)$ & $c_1(q)$\\
\midrule
0&$[0,\frac34]$&$2q+\frac{q^2}2$&$-18-3q$\\
0&$[\frac34,\frac67]$&$-2-\frac23q+\frac{q^2}2$&$-2-3q$\\
0&$[\frac67,1]$&$-4-3q+\frac{q^2}2$&$12-3q$\\
1&$[0,\frac12]$&$-\frac12+2q+\frac{q^2}2$&$-15-3q$\\
1&$[\frac12,\frac57]$&$-\frac{11}6-\frac23q+\frac{q^2}2$&$1-3q$\\
1&$[\frac57,1]$&$-\frac72-3q+\frac{q^2}2$&$15-3q$\\
2&$[0,\frac14]$&$\frac52q+\frac{q^2}2$&$-15-3q$\\
2&$[\frac14,\frac25]$&$-\frac13+\frac76q+\frac{q^2}2$&$-7-3q$\\
2&$[\frac25,\frac47]$&$-\frac23+\frac13q+\frac{q^2}2$&$-2-3q$\\
2&$[\frac47,\frac23]$&$-2-2q+\frac{q^2}2$&$12-3q$\\
2&$[\frac23,1]$&$-3-\frac72q+\frac{q^2}2$&$21-3q$\\
3&$[0,\frac13]$&$1+\frac92q+\frac{q^2}2$&$-33-3q$\\
3&$[\frac13,\frac37]$&$\frac12+3q+\frac{q^2}2$&$-24-3q$\\
3&$[\frac37,\frac35]$&$-\frac12+\frac23q+\frac{q^2}2$&$-10-3q$\\
3&$[\frac35,\frac34]$&$-1-\frac16q+\frac{q^2}2$&$-5-3q$\\
3&$[\frac34,1]$&$-2-\frac32q+\frac{q^2}2$&$3-3q$\\
4&$[0,\frac27]$&$4q+\frac{q^2}2$&$-27-3q$\\
4&$[\frac27,\frac12]$&$-\frac23+\frac53q+\frac{q^2}2$&$-13-3q$\\
4&$[\frac12,1]$&$-2-q+\frac{q^2}2$&$3-3q$\\
5&$[0,\frac17]$&$-\frac12+4q+\frac{q^2}2$&$-24-3q$\\
5&$[\frac17,\frac14]$&$-\frac56+\frac53q+\frac{q^2}2$&$-10-3q$\\
5&$[\frac14,1]$&$-\frac32-q+\frac{q^2}2$&$6-3q$\\
\bottomrule
\end{tabular}
\caption{The law on $22$ pieces: $r=q\bmod6$, the range of $t=\fr{6n/p}$, and the coefficients of
$p\,c_0(q)+n\,c_1(q)$.}\label{tab:pieces}
\end{table}

\begin{proof}
Write $r_3=q\bmod3$, so that $\fr{2n/p}=(r_3+t)/3$, and $\theta=\fr{3n/p}=((q\bmod2)+t)/2$. Up to an error of at most
$2$ each, the class sizes are $p$ times the masses
\[
 m_A=t,\ m_B=0,\ m_C=1-\tfrac{r_3+t}3-t,\ m_D=\tfrac{r_3+t}3\quad\text{if }t<\tfrac{3-r_3}4,
\]
\[
 m_A=1-\tfrac{r_3+t}3,\ m_B=t-m_A,\ m_C=0,\ m_D=1-t\quad\text{otherwise.}
\]
Put $\lambda_c=q+r_3-3$. Since $q_2=(q-r_3)/3$, \eqref{eq:lemT} reads $2d+|\mathcal N|+2=kp$ with
$k=\lambda_c+1-\lambda\ge1$, and $k\le2$ because $d\le|Z|+|W|$. So $\lambda=\lambda_c$, or $\lambda=\lambda_c-1$ and then
$|\mathcal N|\le2$. Let $h(\lambda',\sigma)=j^2+(\sigma-5)j$ with $j=\lceil(\lambda'+4-\sigma)/2\rceil$. This is the
sum of the increments below $\lambda'$ at a residue of class $\sigma$. Let $N_c$ be the number of residues in the two
classes with $\sigma\not\equiv\lambda_c\pmod2$. If $\lambda=\lambda_c$, the sum of the $3n$ smallest increments is
$\sum_c|c|\,h(\lambda_c,\sigma_c)+d\lambda_c$ with $2d=p-N_c-2$. If $\lambda=\lambda_c-1$, then
$h(\lambda_c-1,\sigma)=h(\lambda_c,\sigma)-(\lambda_c-1)[\sigma\not\equiv\lambda_c]$ and $N_c=p-|\mathcal N|$, and the sum
is $\sum_c|c|\,h(\lambda_c,\sigma_c)+\frac12(p-N_c-2)(\lambda_c-1)$. In both cases it equals
$\sum_c|c|\,h(\lambda_c,\sigma_c)+\frac12(p-N_c-2)\lambda_c+O(q)$, and $|h|\le12q^2$. The tie classes for $\lambda_c$ are
$(A,C)$ if $\lfloor q/3\rfloor$ is odd and $(B,D)$ if it is even. Write $m_Z\le m_W$ for their masses and $m_{\mathcal N}$
for the mass of the other two. If $\lambda=\lambda_c$, then $2d+2=|Z|+|W|$, so $x_1=(|W|-|Z|)/2-1$ and
$x_2=(x_1-|\mathcal N|)_+$. If $\lambda=\lambda_c-1$, the classes of the parity of $\lambda_c$ have at most two residues,
so $m_Z$ and $m_W$ are $O(1/p)$. In both cases
\[
 x_1+x_2\ge p\Bigl(\frac{m_W-m_Z}2+\Bigl(\frac{m_W-m_Z}2-m_{\mathcal N}\Bigr)_+\Bigr)-5 .
\]
Also $\nu_p=2\sum_{i<3n}\lfloor i/p\rfloor$ since $p^2>3n$, and this is exactly $p(9x^2-3x+\theta(1-\theta))$ with $x=n/p$. On each piece the result is $p\,c_0(q)+n\,c_1(q)$: the terms in $t^2$
cancel between $\nu_p$ and the increments, and with $t=6x-q$ the rest is affine in $x$. The errors are $O(q^2)$, with
constant at most $1000$.
\end{proof}

\begin{remark}
The formal proof carries out the case analysis of Proposition~\ref{prop:bridge} piece by piece. The constant $1000$
is generous and plays no role, because $\sum_{p>n/M}q^2\log p=O(M^2n)$.
\end{remark}

\subsection{The integrals}
Put $y=p/n$. On a piece, $p\,c_0(q)+n\,c_1(q)=n\bigl(c_0(q)y+c_1(q)\bigr)$, and the piece is the interval
$6/(q+t_1)<y\le6/(q+t_0)$. The weighted prime number theorem $\sum_{\alpha n<p\le\beta n}\log p=(\beta-\alpha+o(1))n$
gives, for each piece $\kappa$ with $q\ge2$ fixed,
\[
 \sum_{p\in\kappa}\bigl(-w(n,p)\bigr)\log p\le\bigl(I_\kappa+o(1)\bigr)n^2,\qquad
 I_\kappa=-\int_\kappa\bigl(c_0(q)y+c_1(q)\bigr)\,dy .
\]

\begin{lemma}\label{lem:bands}
\begin{enumerate}[label=(\roman*)]
\item The pieces with $2\le q\le5$, which cover $n<p<3n$, have $\sum I_\kappa=2263/280$.
\item For $3n<p\le7n/2$ the bound of Theorem~\ref{thm:O} contributes $15/4$ per $n^2$, and for $7n/2<p\le6n$ the bound of
 Theorem~\ref{thm:B7} contributes $47/4$.
\item For $m\ge1$, the pieces with $6m\le q<6m+6$, which cover $n/(m+1)<p\le n/m$, have
 $\sum I_\kappa=f(m):=-\sum_ed_e/(m+e)$, with the $23$ nodes $e$ and weights $d_e$ of Table~\ref{tab:nodes}.
\end{enumerate}
\end{lemma}

\begin{table}[ht]
\centering
\renewcommand{\arraystretch}{1.15}
\setlength{\tabcolsep}{4pt}
\begin{tabular}{@{}c|cccccccccccc@{}}
\toprule
$e$ & $0$ & $\frac18$ & $\frac17$ & $\frac16$ & $\frac14$ & $\frac27$ & $\frac13$ & $\frac38$ & $\frac25$ & $\frac37$ & $\frac49$ & $\frac12$\\
$d_e$ & $-12$ & $8$ & $7$ & $-12$ & $8$ & $7$ & $-\frac{27}2$ & $4$ & $\frac52$ & $7$ & $\frac92$ & $-30$\\
\midrule
$e$ & $\frac59$ & $\frac47$ & $\frac35$ & $\frac58$ & $\frac23$ & $\frac57$ & $\frac34$ & $\frac56$ & $\frac67$ & $\frac78$ & $1$ &\\
$d_e$ & $\frac92$ & $7$ & $\frac52$ & $4$ & $-\frac{27}2$ & $7$ & $8$ & $-12$ & $7$ & $8$ & $-3$ &\\
\bottomrule
\end{tabular}
\caption{The nodes $e$ and weights $d_e$. One has $\sum_ed_e=0$ and $\sum_ed_ee=9/2$.}\label{tab:nodes}
\end{table}

\begin{proof}
(i) and (iii) are exact integrations of linear functions over the pieces. In (iii), on the piece with endpoints
$x_e=m+e$ in the variable $x=1/y$, the substitution $q=6x_e-t_e$ turns the endpoint terms of
$\int(c_0/x^3+c_1/x^2)\,dx$ into a constant, which cancels inside the piece, plus $d/x_e+f/x_e^2$. Summing over the six
residues $r$, the coefficients of $1/x_e^2$ cancel and those of $1/x_e$ are the $d_e$. (ii) is the integral of
$6(y-2)$ over $(3,7/2]$ and of $2+2y$, $18-2y$ and $54-10y$ over $(7/2,4]$, $(4,9/2]$ and $(9/2,6]$.
\end{proof}

Thus the exact part of the content is $2263/280+15/4+47/4=6603/280$. It remains to bound the partial sums of $f$.

\begin{lemma}\label{lem:SM}
Put $S_M=\sum_{1\le m<M}f(m)$ and $R_C=3573565641/700000000$. Then $S_M\le R_C$ for every $M\ge1$.
\end{lemma}

\begin{proof}
\emph{Positivity.} From $e/(m+e)=e/m-e^2/(m(m+e))$ and $\sum d_e=0$, $\sum d_ee=9/2$,
\[
 f(m)=\frac{9/2}{m^2}-\frac1{m^2}\sum_e\frac{d_ee^2}{m+e}\ge\frac9{2m^2}-\frac{\Sigma_2}{m^3},\qquad \Sigma_2=\sum_e|d_e|e^2=\frac{20383}{360}.
\]
So $f(m)>0$ for $m\ge13$, and an exact computation gives $f(m)>0$ for $m\le12$. Hence $S_M$ increases with $M$.

\emph{The tail.} For $K=6$, $1/(m+e)=\sum_{j<K}(-e)^j/m^{j+1}+(-e)^K/(m^K(m+e))$, so
$f(m)=\sum_{k=2}^6a_km^{-k}+r(m)$ with $a_k=(-1)^k\sum_ed_ee^{k-1}$ and $|r(m)|\le\Sigma_6m^{-7}$, $\Sigma_6=\sum_e|d_e|e^6$. Here
$a_2,a_4,a_6>0$ and $a_3,a_5<0$. For $k\ge2$ and $M\ge1$, convexity gives
\[
 \frac1{(k-1)M^{k-1}}+\frac1{2M^k}\le\sum_{m\ge M}m^{-k}\le\frac1{(k-1)(M-\frac12)^{k-1}} .
\]
Using the upper bound for positive coefficients and for $\Sigma_6$, and the lower bound for negative coefficients, we get
$\sum_{m\ge1000}f(m)\le T$ for an explicit rational $T=0.0045006\ldots$.

\emph{The partial sum.} $S_{1000}$ is a sum of $22977$ rationals. Rounding each term $1/(m+e)$ outward to a multiple
of $2^{-80}$ gives an upper bound $\hat S$ with $\hat S-S_{1000}<2\cdot10^{-19}$, and $\hat S+T\le R_C$ with a margin of
$5.76\cdot10^{-9}$. Since $f>0$, $S_M\le S_{1000}+T\le R_C$ for every $M$.
\end{proof}

\begin{theorem}\label{thm:content}
For every $\varepsilon>0$ and all large $n$, $m_n$ is a positive rational number, $m_nJ_n(a,b,X)\in\Z[a,b,X]$, and
\[
 \log m_n\le(\Lambda+\varepsilon)n^2,\qquad\Lambda=\frac{2868723663}{10^8}=28.68723663 .
\]
\end{theorem}

\begin{proof}
Fix $M\ge48/\varepsilon$. Split the primes $p\le6n$ into $p\le n/M$, $n/M<p<3n$ and $3n<p\le6n$. For $p\le n/M$,
Proposition~\ref{prop:crude} gives $\sum_{p\le n/M}12n\lfloor\log_p8n\rfloor\log p\le(12/M+o(1))n^2$. For
$n/M<p<3n$ we use Proposition~\ref{prop:bridge} on each piece. The pieces with $q<6M$ are finite in number, and the error
terms $1000(q^2+1)$ contribute $O(M^2n\log n)$. On $\sqrt{8n}<p\le n/M$ the crude bound is used, so the pieces with
$q\ge6M$ are not needed. Lemma~\ref{lem:bands} gives
\[
 \log m_n=\sum_{p\le6n}(-w(n,p))\log p\le\Bigl(\frac{6603}{280}+S_M+\frac{12}M+\frac\varepsilon2\Bigr)n^2
\]
for large $n$. By Lemma~\ref{lem:SM}, $6603/280+S_M\le6603/280+R_C=2868723663/10^8$.
\end{proof}

The rational bound of Theorem~\ref{thm:content} rounds up an exact constant.

\begin{theorem}\label{thm:C01}
The law on $(0,1]$ has the integral $\sum_{m\ge1}f(m)=\sum_ed_e\psi(1+e)=5.10509376672\ldots$, and
\begin{multline*}
 \frac{6603}{280}+\sum_ed_e\psi(1+e)=95-36\log2+\tfrac{153}2\log3-\tfrac{25}4\log5-49\log7\\
 -8\sqrt2\log(1+\sqrt2)+\tfrac{5\sqrt5}2\log\phi+\tfrac92\Xi,
\end{multline*}
where $\phi=(1+\sqrt5)/2$ and $\Xi=\psi(4/9)+\psi(5/9)+2\gamma$. The value is $28.6872366238631\ldots$.
\end{theorem}

\begin{proof}
Since $\sum_ed_e=0$, $\sum_{m\ge1}f(m)=\sum_ed_e\sum_{m\ge1}\bigl(1/m-1/(m+e)\bigr)=\sum_ed_e(\psi(1+e)+\gamma)=\sum_ed_e\psi(1+e)$.
With $\psi(1+e)=\psi(e)+1/e$ the rational part is $\sum_{e>0}d_e/e=19997/280$. Gauss's digamma theorem evaluates the
symmetric pairs: $\psi(\frac12)=-\gamma-2\log2$, $\psi(\frac13)+\psi(\frac23)=-2\gamma-3\log3$,
$\psi(\frac14)+\psi(\frac34)=-2\gamma-6\log2$, $\psi(\frac16)+\psi(\frac56)=-2\gamma-4\log2-3\log3$,
$\sum_{k=1}^6\psi(\frac k7)=-6\gamma-7\log7$, $\psi(\frac18)+\psi(\frac78)=-2\gamma-8\log2-2\sqrt2\log(1+\sqrt2)$,
$\psi(\frac38)+\psi(\frac58)=-2\gamma-8\log2+2\sqrt2\log(1+\sqrt2)$ and
$\psi(\frac25)+\psi(\frac35)=-2\gamma-\frac52\log5+\sqrt5\log\phi$. The constant $\gamma$ cancels because $\sum d_e=0$.
\end{proof}

\begin{remark}\label{rem:C01}
Theorem~\ref{thm:content} is the form proved in Lean. Its constant $28.68723663$ is the closed form of
Theorem~\ref{thm:C01} rounded up by less than $10^{-8}$. The formal proof checks only the rational partial sums of
Lemma~\ref{lem:SM}, and it does not use Theorem~\ref{thm:C01}. The paper states the closed form because it shows
where the constant comes from. Theorem~\ref{thm:C01} is not part of the formal proof.
\end{remark}

%% =====================================================================
\section{Nonvanishing at prime-tied sizes}\label{sec:nv}

Let $\ell$ be a prime. We show that $J_n(a,b,X)\ne0$ for three kinds of sizes $n$ tied to $\ell$, at all points with
$\ell$-adically bounded coordinates. In all cases the $\ell$-adic valuation of a Gram determinant is computed exactly
by Lemma~\ref{lem:halfint} and its equality case. The first two kinds give one size for each prime. The third gives
every size in an interval. We write $v$ for $v_\ell$ and $\Z_{(\ell)}$ for the rationals with
$v\ge0$.

\subsection{Lai's primes}
\begin{theorem}\label{thm:NV}
Let $m\ge1$, $n=2m$, and let $\ell=9m+1$ be prime. Let $a',b,X'\in\Z_{(\ell)}$ with $3a'+2b\not\equiv0\pmod\ell$. Then
$J_n(a'/\ell,b,X'/\ell^3)\ne0$.
\end{theorem}

For $a'=\ell a$ and $X'=\ell^3c$ with integers $a$, $b$ and $c$ this is Lai's Proposition~3.1 \cite{Lai2026}, where
$b$ is an $\ell$-adic unit.
The new points are the coefficient $a'/\ell$ and the tolerance $v(X)\ge-3$.

\begin{proof}
Write $a=a'/\ell$, $X=X'/\ell^3$ and $\epsilon=3a'+2b$. Put $s=3m=6n-\ell+1$, so that the nodes $k\in[\ell,6n]$ are
$k=\ell+j$ with $0\le j<s$. Every index in the harmonic sums is at most $8n<2\ell$, and $3n<\ell$.

\emph{Step 1: local sizes.} For $k<\ell$, $H_k(i)\in\Z_{(\ell)}$, and for $\ell\le k<2\ell$, $H_k(i)-\ell^{-i}\in\Z_{(\ell)}$. So
$d_k\in\ell^{-1}\Z_{(\ell)}$ for every $k$. For $k<\ell$, $U_k\in\ell^{-1}\Z_{(\ell)}$ and $V_k\in\ell^{-3}\Z_{(\ell)}$. For
$k=\ell+j$, $U_k\equiv-\ell^{-3}\epsilon\pmod{\ell^{-1}}$ and $V_k\in\ell^{-4}\Z_{(\ell)}$. By Kummer's theorem $c_k$ is a unit
for $\ell\le k\le6n$, $\ell^2\mid c_k$ for $s\le k<\ell$, and $c_k\in\Z_{(\ell)}$ for $k<s$.

\emph{Step 2: Lai's lattice.} For $i<s$ let $F_i(t)=\prod_{j<s,j\ne i}(t+j)/(j-i)$ and $\gamma_i=F_i'(-i)$, and put
\[
 P_i=\ell^2\bigl(1-2\gamma_i(t+i)\bigr)F_i^2,\qquad Q_i=\ell(t+i)F_i^2 .
\]
They have degree at most $2s-1=3n-1$, and at $-j$ ($j<s$) one has $P_i(-j)=\ell^2\delta_{ij}$, $P_i'(-j)=0$,
$Q_i(-j)=0$ and $Q_i'(-j)=\ell\delta_{ij}$. The Hermite map $f\mapsto(f(-j),f'(-j))_{j<s}$ is an isomorphism of
$\Z_{(\ell)}[t]_{<2s}$ onto $\Z_{(\ell)}^{2s}$, since its determinant is a power of a Vandermonde determinant in
$0,\dots,s-1$. So the $P_i,Q_i$ span a lattice of index $\ell^{3s}=\ell^{9n/2}$, and
$v\det[B(t^i,t^j)]=v\det G-9n$, where $G$ is the Gram matrix of the $P_i,Q_i$.

\emph{Step 3: the Gram matrix modulo $\ell$.} At the nodes $k<s$, $Q_j(-k)=0$, so every contribution lies in
$\ell\Z_{(\ell)}$. This exact vanishing is what allows $v(X)\ge-3$. At the nodes $s\le k<\ell$, $\ell^2\mid c_k$ and every
contribution lies in $\ell\Z_{(\ell)}$. At $k=\ell+j$, Taylor expansion from $-j$ gives $P_i(-k)\equiv\ell^2\delta_{ij}$,
$Q_i(-k)\equiv-\ell^2\delta_{ij}$ and $Q_i'(-k)\equiv\ell\delta_{ij}$ to the relevant order. With $A_j=c_k\ell^4V_k$ and
$\bar c_j$ the residue of $c_k$, the node contributes, modulo $\ell$,
\[
 \begin{pmatrix}A_j&-\bar c_j\epsilon-A_j\\-\bar c_j\epsilon-A_j&2\bar c_j\epsilon+A_j\end{pmatrix}
\]
to the block of $(P_j,Q_j)$, and nothing else. Its determinant is $-\bar c_j^2\epsilon^2$, whatever $A_j$ is. So
$\det G\equiv\prod_{j<s}(-\bar c_j^2\epsilon^2)\not\equiv0$. Hence $v(\prod i!^2J_n)=-9n$, and $\prod_{i<3n}i!^2$ is a unit.
\end{proof}

The proof gives $v_\ell(J_n(a'/\ell,b,X'/\ell^3))=-9n$.

\begin{corollary}\label{cor:NV}
Let $m\ge1$, $n=2m$, and let $\ell=9m+1$ be prime.
\begin{enumerate}[label=(\alph*)]
\item Let $a,b\in\Z$ with $\ell\nmid b$. Then $J_n(a,b,p/q)\ne0$ for all integers $p$ and $q>0$ with $\ell^4\nmid q$.
\item Let $p,q$ be integers with $q\ne0$ and $\ell\nmid p$. Then $f_n(p/q)\ne0$.
\end{enumerate}
\end{corollary}

\begin{proof}
(a) Take $a'=\ell a$ and $X'=\ell^3p/q$. Then $v(X')\ge0$ and $3a'+2b\equiv2b\not\equiv0$. (b) By \eqref{eq:scaled},
$q^{3n}f_n(p/q)=J_n(q,-p,0)$. Take $a'=\ell q$, $b=-p$ and $X'=0$, so that $3a'+2b\equiv-2p\not\equiv0$.
\end{proof}

\subsection{The lines \texorpdfstring{$3\ell=7n+1$}{3l = 7n+1} and \texorpdfstring{$3\ell=7n+2$}{3l = 7n+2}}
Lai's primes have density $1/6$ among the primes, in the sense of Theorem~\ref{thm:criterion}, and the sizes tied to
them grow with slope $9/2$. A second family of exact nonvanishing comes from primes $\ell$ with $2n<\ell<8n/3$.

\begin{theorem}\label{thm:NV7}
Let $n\ge7$ and let $\ell$ be a prime with $3\ell\in\{7n+1,7n+2\}$ and $\ell\notin\{73,1039\}$. Let $a\in\Z_{(\ell)}$,
$b\in\Z_{(\ell)}^\times$ and $X'\in\Z_{(\ell)}$. Then $J_n(a,b,X'/\ell^3)\ne0$.
\end{theorem}

\begin{proof}
Put $\alpha_0=\ell-2n$ and $A_2=6n-2\ell$, and write $k=s+j\ell$ with $0\le s<\ell$. Since $6n<3\ell$, $j\le2$. In the
notation of Section~\ref{sec:local} the residues fall into the classes $T_1=[0,\alpha_0)$ with $\sigma=0$ and $j\le2$,
$T_2=[\alpha_0,A_2]$ with $\sigma=3$ and $j\le2$, and $T_4=(A_2,\ell)$ with $\sigma=5$ and $j\le1$. Put $\varkappa=3\ell-7n$.
Then $|T_2|=n+1-\varkappa$ and $|T_4|=n-1+\varkappa$.

\emph{Step 1: node sizes.} For $j\ge1$, Lemma~\ref{lem:N} gives $U_k=\ell^{-3}u_jb+O(\ell^{-2})$ and
$V_k=\ell^{-4}E_jb+O(\ell^{-3})$, with $u_1=-2$, $u_2=-9/4$, and $E_j=-6,-105/8$ on $T_1$, $E_j=-3,-87/8$ on $T_2$ and
$E_1=-7$ on $T_4$. The terms in $a$ and $X$ are one order smaller. For $j=0$, $v(U_k)\ge0$ and $v(V_k)\ge-3$. The unit parts
of $c_{s+j\ell}$ are $(1,4,1)\hat c_s$ on $T_1$ and $(1,32,27)\hat c_s$ on $T_2$.

\emph{Step 2: the basis.} Put $\Pi=\prod_{r\in T_1}(t+r)^3\prod_{r\in T_2}(t+r)$, of degree $2n+1$, and
$\Psi_r=\Pi/(t+r)^3$. The basis consists of $e_{r,i}=(t+r)^i\Psi_r$ for $r\in T_1$, $i<3$, with potential $i-2$; of
$e_r=\Pi/(t+r)$ for $r\in T_2$, with potential $-\frac12$; and of $g_m=t^m\Pi$ for $m<n-1$, with potential $\frac12$. It
has $3n$ elements and $2\sum\pi=5n-3\ell-2$. By Lemma~\ref{lem:N} and Step 1 every entry has $v\ge\pi_x+\pi_y$, with
$+\frac12$ between the two groups. At the nodes with $j=0$ every term lies strictly above its threshold: on $T_1$ only
$e_{s,0}$ and $e_{s,1}$ do not vanish to order $3$, on $T_2$ only $e_s$ does not vanish, and on $T_4$ the weight has
$v=5$. So $X$ never enters the reductions. In the proof of Lemma~\ref{lem:halfint} the reduction is block triangular, so
\[
 v(\det G)=2\sum\pi\iff\det\overline N_{AA}\cdot\det\overline N_{BB}\ne0\ \text{in }\F_\ell .
\]
The reductions use only the leading terms, so they are $\bar b$ times their values at $(a,b,X)=(0,1,0)$.

\emph{Step 3: $\overline N_{AA}$.} The elements of different residues of $T_1$ meet only above the threshold. At the
nodes $-r-j\ell$, $e_{r,i}=(-j\ell)^i\Psi_r(-r)+O(\ell^{i+1})$. Hence $\overline N_{AA}$ is block diagonal with blocks
$\hat c_r\Psi_r(-r)^2M_0$, where $M_0=[h_{i+i'}]_{i,i'<3}$ and $h_m=\sum_{j=1,2}\binom2j^2(-j)^{m-1}(mu_j-jE_j)$. One finds
$h=(-\frac{297}8,40,-\frac{103}2,78,-130)$ and $\det M_0=1095/2=3\cdot5\cdot73/2$.

\emph{Step 4: $\overline N_{BB}$.} The element $e_r$ meets $e_{r'}$ only for $r=r'$, at the nodes of $T_2$, with value
$-\frac{3117}8\hat c_r\Pi'(-r)^2$. The pairing of $e_r$ with $g_m$ is $\frac{1117}2\hat c_r\Pi'(-r)^2(-r)^m$. The pairing of
$g_m$ with $g_{m'}$ is $-\frac{1799}2\hat c_r\Pi'(-r)^2(-r)^{m+m'}$ at each residue of $T_2$ plus
$-7\hat c_{s,1}\Pi(-s)^2(-s)^{m+m'}$ at each residue of $T_4$. By Anton's congruence \cite[Theorem~1]{Granville97} and Wilson's theorem,
$\hat c_s\Pi'(-s)^2=-KD(s)$ on $T_2$ and $\hat c_{s,1}\Pi(-s)^2=-32KD(s)$ on $T_4$, with $D(s)=\prod_{i<\alpha_0}(s-i)$ and a unit
$K$. The diagonal block is invertible since $3117=3\cdot1039$, and its Schur complement is
\[
 S_{mm'}=-K\Bigl[c_2\sum_{T_2}D(r)(-r)^{m+m'}-224\sum_{T_4}D(s)(-s)^{m+m'}\Bigr],
\]
where
\[
 c_2=-\frac{1799}2+\Bigl(\frac{1117}2\Bigr)^2\frac8{3117}=-\frac{616727}{6234}.
\]
The polynomial $D(s)s^{m+m'}$ vanishes on $T_1$ and has degree at most $\ell-4$, so its sums over $T_2$ and $T_4$ are
opposite. Hence $S=K\frac{779689}{6234}[\sum_{s\in T_4}D(s)(-s)^{m+m'}]$, with $779689=883^2$, and it remains to show
\[
 \mathcal H=\det\Bigl[\sum_{s\in T_4}D(s)s^{m+m'}\Bigr]_{m,m'<n-1}\ne0 .
\]

\emph{Step 5: duality.} For points $x_s$ with weights $\omega_s$ ($s=1,\dots,N$) and $\nu<N$, write
$H_{\nu}(\omega)=\det[\sum_s\omega_sx_s^{i+j}]_{i,j<\nu}$. The Cauchy--Binet formula expresses it as a
sum over subsets of size $\nu$ of $\prod\omega$ times a squared Vandermonde determinant. Passing to the
complements gives
\[
 H_{\nu}(\omega)=\prod_s\omega_s\cdot V^2\cdot H_{N-\nu}(\omega^*),\qquad
 \omega^*_s=\frac1{\omega_s\prod_{t\ne s}(x_s-x_t)^2},
\]
with $V$ the Vandermonde determinant of all the points. We apply this to $T_4$ with $\omega=D$, $\nu=n-1$ and
$N-\nu=\varkappa$. By Wilson's theorem, $s!(\ell-1-s)!\equiv(-1)^{s+1}$, so
$\prod_{t\in T_4,t\ne s}(s-t)^2\equiv\bigl((s-A_2-1)!/s!\bigr)^2$. Since $\alpha_0\le A_2+1$, the dual weight is the
polynomial
\[
 \omega^*_s\equiv\Phi(s)=\prod_{i\le A_2}(s-i)\cdot\prod_{i=\alpha_0}^{A_2}(s-i),
\]
which is monic of degree $\ell+2-2\varkappa$ and vanishes off $T_4$. So the dual entries are full power sums
$\tau_m=\sum_{s\in\F_\ell}\Phi(s)s^m$, and $\sum_{s\in\F_\ell}s^j=-1$ if $(\ell-1)\mid j>0$ and $0$ otherwise. For
$\varkappa=1$, $H_1=\tau_0=-[s^{\ell-1}]\Phi$ is the sum of the roots of $\Phi$, which is $(31n^2-7n-2)/18\equiv-1/49$,
since $n\equiv-1/7$. For $\varkappa=2$, $\tau_0=0$, $\tau_1=-1$, and $H_2=\tau_0\tau_2-\tau_1^2=-1$.

\emph{Step 6.} On both lines $\det\overline N_{AA}\det\overline N_{BB}\ne0$ unless $\ell$ divides
$2\cdot3\cdot5\cdot7\cdot73\cdot883\cdot1039$. Since $n\ge7$ gives $\ell\ge17$, and $3\cdot883$ lies on neither line, this
excludes only $\ell=73$ (with $n=31$) and $\ell=1039$ (with $n=445$).
\end{proof}

The proof gives $v_\ell(J_n(a,b,X'/\ell^3))=5n-3\ell-2-2(3n-\ell)=-n-\ell-2$.

\begin{corollary}\label{cor:NV7}
Let $n$ and $\ell$ be as in Theorem~\textup{\ref{thm:NV7}}. If $a,b\in\Z$ with $\ell\nmid b$, then $J_n(a,b,p/q)\ne0$
whenever $\ell^4\nmid q$. If $\ell\nmid p$ and $q\ne0$, then $f_n(p/q)\ne0$.
\end{corollary}

\begin{proof}
As for Corollary~\ref{cor:NV}, with $(a,b,X')=(a,b,\ell^3p/q)$ and $(q,-p,0)$.
\end{proof}

\subsection{The band \texorpdfstring{$8n/3<\ell<3n-1$}{8n/3 < l < 3n-1}}\label{ssec:band}
The two families above give nonvanishing at one size $n$ for each prime. In this subsection a prime $\ell$ is good for
every size $n$ with $8n/3<\ell<3n-1$, and for every direction. The proof shows that a reduced Gram matrix is
nonsingular modulo $\ell$.

Let $n$ be a positive integer and let $\ell$ be a prime with $8n<3\ell$ and $\ell+1<3n$. Then $8n<3\ell<9n-3$, so
$n\ge4$. No integer $\ell$ satisfies both inequalities for $n=4$, $5$ or $6$, so $n\ge7$ and $\ell>8n/3$ gives
$\ell\ge19$. Put
\[
 \alpha_0=\ell-2n,\qquad A_2=6n-2\ell,\qquad \gamma=\frac{\alpha_0-1}2,\qquad g=3n-1-\ell .
\]
Then $\alpha_0$ is odd, $6n=2\ell+A_2$, $0\le A_2<\alpha_0$, $A_2=2g+2$ and
\begin{equation}\label{eq:band}
 \ell=3\alpha_0+A_2=6\gamma+2g+5 .
\end{equation}
In the notation of Section~\ref{sec:local} we have $q_6=2$, $r_6=A_2$, $q_2=0$ and $r_2=2n$, so the class $B$ is empty. The
residues $s\in[0,\ell)$ form three intervals:
\[
 \begin{aligned}
 A&=[0,A_2]&&(\sigma=0,\ k=s+j\ell,\ j\le2),\\
 C&=(A_2,\alpha_0)&&(\sigma=2,\ j\le1),\\
 D&=[\alpha_0,\ell)&&(\sigma=5,\ j\le1).
 \end{aligned}
\]
We put $\mathrm{w}_{n,\ell}=-21n+7\ell-7$.

\begin{theorem}\label{thm:B83}
Let $n$ be a positive integer and let $\ell$ be a prime with $8n/3<\ell<3n-1$ and $\ell\notin\{2,3,5,13,24359\}$.
Let $a\in\Z_{(\ell)}$, $b\in\Z_{(\ell)}^\times$ and $X\in\Q$ with $\ell^3X\in\Z_{(\ell)}$. Then $J_n(a,b,X)\ne0$.
\end{theorem}

Since the band forces $\ell\ge19$, the exclusion of $2$, $3$, $5$ and $13$ is automatic. The prime $24359$ lies in the
band exactly for $8120<n\le9134$. The proof gives more than nonvanishing. It shows that
$v_\ell(J_n(a,b,X))=\mathrm{w}_{n,\ell}$, and this value is the bound of Theorem~\ref{thm:L2} in this range. So the
law is sharp at every such point. This exact valuation is not part of the formal proof, which uses only the
nonvanishing.

\begin{corollary}\label{cor:B83}
Let $n$ and $\ell$ be as in Theorem~\ref{thm:B83}.
\begin{enumerate}[label=(\alph*)]
\item Let $a,b\in\Z$ with $\ell\nmid b$. Then $J_n(a,b,p/q)\ne0$ for all integers $p$ and $q>0$ with $\ell^4\nmid q$.
\item Let $p,q$ be integers with $q\ne0$ and $\ell\nmid p$. Then $f_n(p/q)\ne0$.
\end{enumerate}
\end{corollary}

\begin{proof}
(a) Apply Theorem~\ref{thm:B83} with $X=p/q$. Since $\ell^4\nmid q$, we have $\ell^3X\in\Z_{(\ell)}$. (b) By \eqref{eq:scaled}, $q^{3n}f_n(p/q)=J_n(q,-p,0)$, and
Theorem~\ref{thm:B83} applies with $(a,b,X)=(q,-p,0)$.
\end{proof}

The proof has four parts: a basis and a reduction to one point, the integral blocks, the half-integral block and its
Hankel determinant, and the assembly.

\subsubsection{The basis and the reduction to one point}
Let $\mu_r=4$ for $r\in A$, $\mu_r=3$ for $r\in C$, and $\Pi=\prod_{r<\alpha_0}(t+r)^{\mu_r}$. Put
\[
 e_{r,i}=(t+r)^i\,\Pi/(t+r)^{\mu_r}\quad(r<\alpha_0,\ i<\mu_r),\qquad z_j=t^j\Pi\quad(j<g).
\]
The roots $-r$ are pairwise incongruent modulo $\ell$, and $\sum_r\mu_r+g=3\alpha_0+A_2+1+g=\ell+1+g=3n$. So this is a
basis in the sense of Lemma~\ref{lem:gram}. Give $e_{r,i}$ the potential $i-2$ if $r\in A$ and $i-1$ if $r\in C$, and give
$z_j$ the potential $\frac32$. The elements $e_{r,i}$ form the \emph{integral group} and the $z_j$ the
\emph{half group}. Put $K_0=32(\ell-1)!^3A_2!^2/(2n)!^3$, a unit. Then $2\sum\pi=-4(A_2+1)+3g=5\ell-15n-7$. Also $\nu_\ell=2(3n-\ell)$, because $\ell^2>3n$. So
\begin{equation}\label{eq:target}
 2\sum\pi-\nu_\ell=\mathrm{w}_{n,\ell}.
\end{equation}

The entries of the Gram matrix are linear forms in $(a,b,X)$. Write $G_{xy}=aG^a_{xy}+bG^b_{xy}+XG^X_{xy}$, and
let $\mathrm{thr}(x,y)=\pi_x+\pi_y$ if $x,y$ lie in the same group and $\pi_x+\pi_y+\frac12$ otherwise.

\begin{lemma}[entries]\label{lem:entries}
For every pair of basis elements, $v(G^b_{xy})\ge\mathrm{thr}$, $v(G^a_{xy})\ge\mathrm{thr}+1$ and $v(G^X_{xy})\ge\mathrm{thr}+4$.
\end{lemma}

\begin{proof}
We use the elementary fact that $\sum_{s\in\F_\ell}H(s)=0$ for $H\in\F_\ell[s]$ of degree at most $\ell-2$ (M1), and
\eqref{eq:faulhaber} and \eqref{eq:shift} (M2).

\emph{The $b$-part.} In the band $\lambda=1$ is odd, so the tie classes are $(B,D)$ with $Z=B=\emptyset$ and $W=D$, and
$\mu_W=0$. The basis above is the basis of the proof of Theorem~\ref{thm:L2}, which gives $v(G^b)\ge\mathrm{thr}$.

\emph{Own nodes, $j\ge1$.} By Lemma~\ref{lem:N}(ii) the terms in $a$ are smaller than those in $b$ by one power of
$\ell$, and the term in $X$ is $\ell$-integral, so by four powers. Hence the bounds with $+1$ and $+4$ hold at every node of $A$ and
$C$ with $j\ge1$.

\emph{Nodes with $j=0$ on $A$ and $C$.} At these nodes $v(U_k)\ge0$ and $v(V_k)\ge-1$ in each of the three
variables, so the parts in $a$ and $b$ have the same lower bound. As in the proof of Theorem~\ref{thm:L2}, this
bound exceeds the threshold by $3$, which covers the $a$-part. The term in $X$ of the node $k=r$ is $6c_rF(-r)$. Here
$F(-r)$ vanishes unless both elements are $e_{r,0}$. In that case the coefficient of $X$ has $v\ge\sigma(r)$, and
$\sigma(r)=\mathrm{thr}+4$, since $\mathrm{thr}=-4$ on $A$ and $-2$ on $C$.

\emph{The $D$-nodes, $a$-part.} At $k=\ell+s$ with $s\in D$ the leading order of the $a$-part is $\ell^2$, one above the $b$-part. The
thresholds $1$ and $2$ hold outright, or after (M1) when the polynomial $G=\Omega\psi$ of Lemma~\ref{lem:G3} vanishes at the
residues outside $D$. For $z_j\times z_{j'}$ and (top element of $A$)$\times z_j$ the threshold is $3$. The node data of
Lemma~\ref{lem:G3} below give $-\frac{15}2-3\ell\lambda_0$ and $-3$ for the coefficients of
$F$ and $\ell F'$ in the $a$-part. So $\sum_D(\cdots)/\ell^2\equiv-\frac{15}2\ell\int G+\ell\frac{15}2\int G=0\pmod{\ell^2}$ by (M2), because the total weight of the harmonic numbers in
$\lambda_0$ is $0$ and in the expansion of $c_{\ell+s}$ it is $-1$.

\emph{The $D$-nodes, $X$-part.} Here $c_k\sim\ell^5$ for $j=0$ and $j=1$, and $c_{\ell+s}/c_s\equiv8(1+\ell(2H_{s_2}-5H_s+3H_{s_1}))$
with $s_1=s-\alpha_0$ and $s_2=s-A_2-1$. For the pairs of threshold $3$ the sum
$\sum_D[c_sF(-s)+c_{\ell+s}F(-\ell-s)]/(\ell^5K_0/8)$ is $\equiv9\ell\int G+\ell(-1+8(-1))\int G=0\pmod{\ell^2}$ by (M2). So $v\ge7=\mathrm{thr}+4$. The other pairs need at most $6$ and follow from (M1).
\end{proof}

The next lemma moves the computation to one point. 
\begin{lemma}[transfer]\label{lem:PD}
Let $F\in\Z_\ell[\alpha,Y]$ with $F(0,0)\in\Z_\ell^\times$. Then $F(\alpha,Y)\in\Z_\ell^\times$ for all $\alpha,Y\in\ell\Z_\ell$.
\end{lemma}

\begin{proof}
$F(\alpha,Y)\equiv F(0,0)\pmod\ell$.
\end{proof}

We apply this to $F(\alpha,Y)=\ell^{-\mathrm{w}_{n,\ell}}J_n(\alpha/\ell,1,Y/\ell^4)$. By Lemma~\ref{lem:entries} the matrix
$N(\alpha,Y)=D^{-1}GD^{-1}$ of Lemma~\ref{lem:halfint}, with $G=G^b+(\alpha/\ell)G^a+(Y/\ell^4)G^X$, has entries in $\Z_\ell[\varpi][\alpha,Y]$, where $\varpi^2=\ell$. So
$F$ has integral coefficients. For $a\in\Z_{(\ell)}$ and $\ell^3X\in\Z_{(\ell)}$ the point $\alpha=\ell a$, $Y=\ell^4X$
lies in $\ell\Z_\ell$. By the homogeneity \eqref{eq:homog}, $J_n(a,b,X)=b^{3n}J_n(a/b,1,X/b)$ for a unit $b$. Hence
Theorem~\ref{thm:B83} holds as soon as the reduction of $N(0,0)$ has nonzero determinant.

By the proof of Lemma~\ref{lem:halfint} this reduction is block triangular. Its diagonal blocks are
$\overline N_{\mathrm{int}}$ on the integral group and $\overline N_{\mathrm{half}}$ on the half group. Recall that
$\ell\ge19$. We show that both blocks have nonzero determinant when $\ell\ne24359$.
At $(a,b,X)=(0,1,0)$ we have $U_k=-2H_k(3)$ and $V_k=U_kd_k-6H_k(4)$, and $L(F)=\sum_kc_k(U_kF'(-k)+V_kF(-k))$.

\subsubsection{The integral blocks}
Put $\Psi_r=\Pi/(t+r)^{\mu_r}$, and let $\hat c_r$ be the unit part of $c_r$,
so that $c_{r+j\ell}\equiv\kappa_j\hat c_r$ by Lucas's theorem, with $\kappa=(1,4,1)$ on $A$ and $(1,1)$ on $C$. Let
$P_r(s)=D_{\alpha_0}(s)/(s-r)$ with $D_m(s)=\prod_{i<m}(s-i)$.

\begin{lemma}[unit identities]\label{lem:U83}
For $r\in A$, $K_0P_r(r)^3\equiv32\,\hat c_r\Psi_r(-r)^2\pmod\ell$. For $r\in C$, $K_0P_r(r)^3\equiv8\,\hat c_r\Psi_r(-r)^2$.
\end{lemma}

\begin{proof}
By Anton's congruence, $c_{r+\ell}\equiv4\hat c_r$ on $A$, and $c_{r+\ell}/\ell^2\equiv c_r/\ell^2$ on $C$. Write
$P_r(r)=r!(-1)^{\alpha_0-1-r}(\alpha_0-1-r)!$ and $\Psi_r(-r)^2$ as factorials. The ratio
$K_0P_r(r)^3/(\hat c_r\Psi_r(-r)^2)$ reduces to $-32(-1)^{\alpha_0-1-r}/[(2n+r)!(\alpha_0-1-r)!]^3$ on $A$, and to
$-8(-1)^{\alpha_0-1-r}/[\cdots]^3$ on $C$. Since $(2n+r)+(\alpha_0-1-r)=\ell-1$, Wilson's theorem gives
$(2n+r)!(\alpha_0-1-r)!\equiv(-1)^{r+1}$. With $\alpha_0$ odd, this gives $32$ and $8$.
\end{proof}

\begin{lemma}[the blocks of $A$]\label{lem:blockA}
For $r\in A$ the reduced block of $e_{r,0},\dots,e_{r,3}$ is $\hat c_r\Psi_r(-r)^2M_A$, where
$M_A=[\chi_{i+i'}]_{i,i'<4}$ and $\chi=(-\frac{297}8,40,-\frac{103}2,78,-130,0,0)$. Its determinant is $-16466684=-2^2\cdot13^2\cdot24359$.
\end{lemma}

\begin{proof}
\emph{Own nodes.} At $k=r+j\ell$ ($j=1,2$), Lemma~\ref{lem:N} gives $U_k=\ell^{-3}u_j+\cdots$ and $V_k=\ell^{-4}E_j+\cdots$ with
$u=(-2,-\frac94)$, $E=(-6,-\frac{105}8)$. For the pair $(i,i')$ with $m=i+i'$ the own nodes contribute
$\Phi_A(m)\hat c_r\Psi_r(-r)^2\ell^{m-4}$ with
$\Phi_A(m)=\sum_{j=1,2}\kappa_j(-j)^{m-1}(mu_j-jE_j)$. This gives $\Phi_A=(-\frac{297}8,40,-\frac{103}2,78,-130,224,-384)$ for $m=0,\dots,6$.

\emph{The $D$-nodes.} At $k=\ell+s$, $s\in D$, the contribution of the pair is $-7\ell K_0\sum_DG+\ell^2K_0S_1+\cdots$ with
$G=D_{\alpha_0}^3\psi$ and $\psi(s)=(r-s)^{m-8}$. For $m\le4$ it lies above the threshold. For $m=5$, $G$ is regular at $r$ with $G(r)=-P_r(r)^3$, so
$\sum_DG\equiv-G(r)$ by (M1), and the entry is $\ell[224-7\cdot32]\hat c_r\Psi_r(-r)^2=0$ by Lemma~\ref{lem:U83}. For $m=6$,
$G=(s-r)P_r^3$ has a simple zero at $r$. At level $\ell^2$ the $D$-part is
$K_0[5(\Omega\psi')(r)+10(\Omega\Lambda_1\psi)(r)-18G'(r)]=12K_0P_r(r)^3=384\,\hat c_r\Psi_r(-r)^2$. The terms in $\int G$ cancel,
and the shift corrections are $G'(r)$ for the shifts $s-\alpha_0$ and $s-A_2-1$. So the entry is $-384+384=0$.

The determinant of $M_A$ is an exact computation with rational numbers.
\end{proof}

\begin{lemma}[the blocks of $C$]\label{lem:blockC}
For $r\in C$ the reduced block of $e_{r,0},e_{r,1},e_{r,2}$ is $\hat c_r\Psi_r(-r)^2\bigl[\begin{smallmatrix}-10&8&-6\\8&-6&-52\\-6&-52&-130\end{smallmatrix}\bigr]$,
and its determinant is $2^{15}$.
\end{lemma}

\begin{proof}
The own node $j=1$ gives $\Phi_C(m)=(-1)^{m-1}(10-2m)$. For $m=3$ the $D$-part is $-7\ell K_0P_r(r)^3=-56\,\hat c_r\Psi_r(-r)^2\ell$, so the entry is $4-56=-52$.
For $m=4$ the $D$-part is $K_0P_r^3[5(-2)+10\cdot3-36]=-16K_0P_r^3=-128\,\hat c_r\Psi_r(-r)^2\ell^2$, since $\mu_r=3$ and only the
shift $s-\alpha_0$ crosses $r$. The entry is $-2-128=-130$. The determinant is $2t(245t+88)$ at $t=8$, which is $2^{15}$.
\end{proof}

Entries between elements of different residues $r\ne r'$ vanish in the reduction. At the own nodes they lie strictly above the threshold, since
an element that is not own at a node vanishes there to order $\mu_r$. At the $D$-nodes the pairs of threshold $2$ have
$G$ vanishing to order at least $2$ at both residues. Then (M1) and (M2) give $S_1\equiv7\int G$, which cancels $-7\int G$. The pairs of threshold at most $1$ follow from (M1).
Hence
\[
 \det\overline N_{\mathrm{int}}=\prod_{r\in A}(\hat c_r\Psi_r(-r)^2)^4\det M_A\cdot\prod_{r\in C}(\hat c_r\Psi_r(-r)^2)^3\,2^{15},
\]
which is nonzero for $\ell\notin\{2,13,24359\}$.

\subsubsection{The half-integral block}
Let $\mathcal V$ be the linear functional on $\Q[x]$ with $\mathcal V(x^e)=eB_{e-1}$, where $B_1=-\frac12$, and put $\int g=\sum_eg_eB_e$, so
that $\mathcal V(g)=\int g'$. Let $\Omega=D_{\alpha_0}^3$.

\begin{lemma}\label{lem:G3}
Let $\varphi\in\Z_{(\ell)}[t]$ with $\deg\varphi\le2g-2$, and put $\psi(x)=\varphi(-x)$ and $G=\Omega\psi$. Then
\[
 L(\Pi^2\varphi)\equiv15K_0\,\ell^3\,\mathcal V(G)\pmod{\ell^4}.
\]
\end{lemma}

The degree of $G$ is at most $3\alpha_0+2g-2=\ell-4$ by \eqref{eq:band}. The proof has four steps.

\emph{Step 1: summation lemmas.} Let $g_0\in\Z_{(\ell)}[x]$ have degree at most $\ell-3$ and put $\mathcal G(x)=\sum_{u<x}g_0(u)$. This is a
polynomial divisible by $x$, by Faulhaber's formula. Put $q=\mathcal G/x$ and $\mathcal T(g_0)=\int q$. By the theorem of von Staudt, $\int g_0$, $\mathcal V(g_0)$ and $\mathcal T(g_0)$ lie in $\Z_{(\ell)}$.
Wolstenholme's theorem gives $H_{\ell-1}\equiv0\pmod{\ell^2}$ and $H_{\ell-1}(2)\equiv0\pmod\ell$. With Faulhaber's formula this gives
\begin{align*}
 \sum_{u<\ell}g_0(u)&\equiv\ell\!\int g_0+\tfrac{\ell^2}2\mathcal V(g_0)\ (\ell^3),\\
 \sum_{u<\ell}g_0(u)H_u&\equiv\int g_0-\ell\,\mathcal T(g_0)\ (\ell^2),\\
 \sum_{u<\ell}g_0(u)H_u(2)&\equiv\tfrac12\mathcal V(g_0)\ (\ell),\\
 \sum_{u<\ell}g_0(u)H_u^2&\equiv-2\mathcal T(g_0)-\tfrac12\mathcal V(g_0)\ (\ell).
\end{align*}
For the second, $\sum_ug_0(u)H_u=\mathcal G(\ell)H_{\ell-1}-\sum_{i=1}^{\ell-1}q(i)$ by summation by parts, and the first formula applied to $q$ gives the claim, since $q(0)=\int g_0$. The third is the same computation
with $\mathcal G(i)/i^2$. For the fourth, $H_u^2=H_u(2)+2\sum_{j\le u}H_{j-1}/j$, and the second formula for $q$ finishes. By Lemma~\ref{lem:Wt} below these apply to every
weight in the proof, and sums over $D$ equal sums over $[0,\ell)$.

\emph{Step 2: the nodes.} Write $k=j\ell+s$. By Lemma~\ref{lem:N}, $v(c_k)=\sigma(s)$. The polynomial $F=\Pi^2\varphi$ vanishes to order $8$ at $-r$ for $r\in A$ and to order $6$ for $r\in C$, and $H_k(3)$, $H_k(4)$, $d_k$ have $v\ge-3,-4,-1$. So the nodes with $s\in A\cup C$ and $j=0$ vanish. The nodes with $s\in A$ and $j\ge1$ have $v\ge0-4+7$, those with $s\in C$ and $j=1$ have $v\ge2-4+5$, and those with $s\in D$ and $j=0$ have $v\ge5-1$. Hence every node other than $k=\ell+s$ with $s\in D$ contributes $0$ modulo $\ell^4$.
So $L(\Pi^2\varphi)\equiv T=\sum_{s\in D}[c_{\ell+s}V_{\ell+s}F(-\ell-s)+c_{\ell+s}U_{\ell+s}F'(-\ell-s)]\pmod{\ell^4}$.

\emph{Step 3: the node term.} Fix $s\in D$, $k=\ell+s$, $s_1=s-\alpha_0$, $s_2=s-A_2-1$. Write
$\langle x,y\rangle=1+\ell x+\ell^2(y+x^2/2)$, so that $\langle x,y\rangle\langle x',y'\rangle\equiv\langle x+x',y+y'\rangle\pmod{\ell^3}$ and
$\prod_{y=1}^{m}(1+j\ell/y)\equiv\langle jH_m,-\frac{j^2}2H_m(2)\rangle$. Digit expansions of the factorials, the reflection
$(\ell-1)!=(\ell-1-m)!(-1)^mm!\prod_{y\le m}(1-\ell/y)$ and Wolstenholme's theorem give
\[
 c_{\ell+s}\equiv\ell^5K_0'\frac{s_1!^3s_2!^2}{s!^5}\bigl\langle\kappa_1+6H_{s_1}-5H_s-2H_{s_2},\ \kappa_2-6H_{s_1}(2)+\tfrac52H_s(2)-H_{s_2}(2)\bigr\rangle
\]
modulo a relative error $\ell^3$, with $K_0'=32A_2!^2(\alpha_0-1)!^3$ and constants $\kappa_1,\kappa_2$ independent of $s$. Also
$\Pi(-s)^2=s!^8s_2!^{-2}s_1!^{-6}$ exactly, so $s_1!^3s_2!^2\,\Pi(-s)^2/s!^5=\Omega(s)$, and
$\Pi(-k)^2/\Pi(-s)^2\equiv\langle2\Lambda_1(s),-\Lambda_2(s)\rangle$ with $\Lambda_j(s)=\sum_r\mu_r(s-r)^{-j}$. Next, $H_k(3)=\ell^{-3}+(\text{integral})$ and
$H_k(4)=\ell^{-4}+(\text{integral})$, and
\[
 \ell d_k\equiv\tfrac12+\ell\lambda_0+\ell^2\delta_1,\qquad \ell^4V_k\equiv-7-2\ell\lambda_0-2\ell^2\delta_1\pmod{\ell^3},\qquad \ell^3U_k\equiv-2,
\]
with $\lambda_0=5H_s-3H_{s_1}-2H_{s_2}$ and $\delta_1=-5H_s(2)+6H_{s_1}(2)-2H_{s_2}(2)$. With $\psi(x)=\varphi(-x)$ we have
$F(-k)=\Pi(-k)^2\psi(k)$ and $F'(-k)=\Pi(-k)^2(-2\Lambda_1(k)\psi(k)-\psi'(k))$. Taylor expansion gives, modulo $\ell^3$,
\begin{multline*}
 \ell^4\bigl(V_kF(-k)+U_kF'(-k)\bigr)/\Pi(-k)^2\equiv2\ell(\psi'+\ell\psi'')\\
 +(-7-2\ell\lambda_0-2\ell^2\delta_1+4\ell\Lambda_1-4\ell^2\Lambda_2)(\psi+\ell\psi'+\tfrac{\ell^2}2\psi'').
\end{multline*}

\emph{Step 4: collection.} By the exact identities $H_{s_1}=H_s-\delta_N(s)$, $H_{s_2}=H_s-\delta_A(s)$, and the same for the second order sums, with
$\delta_N(x)=\sum_{i<\alpha_0}(x-i)^{-1}$ and $\delta_A(x)=\sum_{i\le A_2}(x-i)^{-1}$, all harmonic numbers collapse to $H_s$, $H_s(2)$ and
rational functions with poles at the zeros of $D_{\alpha_0}$. Since $\Lambda_1=3\delta_N+\delta_A$, we get
$e_1=-H_s+4\delta_A$, $\varepsilon_2=-\frac92H_s(2)+3\delta_N^{(2)}$, $-2\lambda_0+4\Lambda_1=6\delta_N$ and $-2\delta_1-4\Lambda_2=2H_s(2)-8\delta_A^{(2)}$,
where $e_1,\varepsilon_2$ are the arguments of $\langle\cdot,\cdot\rangle$ in the expansion of $c_{\ell+s}\Pi(-k)^2$ beyond the constants. With
$W_1=\Omega\psi'$ and $W_2=\Omega\psi''$, the expansion of $\Omega\langle e_1,\varepsilon_2\rangle\times(\text{bracket})$ is
\begin{multline*}
 -7G+\ell\bigl[-5W_1+7GH_s-28G\delta_A+6G\delta_N\bigr]\\
 +\ell^2\bigl[-\tfrac32W_2-\tfrac72GH_s^2+\tfrac{67}2GH_s(2)+5H_sW_1+28GH_s\delta_A-6GH_s\delta_N+R\bigr],
\end{multline*}
where $R$ is a combination, with integral coefficients, of terms without harmonic numbers. The functions $G$, $W_1$, $W_2$, $G\delta_N$, $G\delta_A$, $G\delta_N^{(2)}$,
$G\delta_A^{(2)}$, $G\delta_A^2$, $G\delta_A\delta_N$, $W_1\delta_A$ and $W_1\delta_N$ are polynomials of degree at most $\ell-4$ that vanish on $[0,\alpha_0)$
(Lemma~\ref{lem:Wt}). Summing the three brackets over $D$ with Step 1, with the accuracies $\ell^3$, $\ell^2$ and $\ell$:
\begin{itemize}
\item bracket $0$: $-7\ell\int G-\frac72\ell^2\mathcal V(G)$;
\item bracket $1$: $7\int G-7\ell\mathcal T(G)-5\ell\int W_1-28\ell\int G\delta_A+6\ell\int G\delta_N$;
\item bracket $2$: $7\mathcal T(G)+\frac74\mathcal V(G)+\frac{67}4\mathcal V(G)+5\int W_1+28\int G\delta_A-6\int G\delta_N$, the terms without harmonic numbers being $0$ by the first formula of Step 1.
\end{itemize}
In the total of bracket $0$, $\ell$ times bracket $1$ and $\ell^2$ times bracket $2$, the coefficient of $\ell$ is $-7\int G+7\int G=0$. The terms $\mathcal T$, $\int W_1$, $\int G\delta_A$ and $\int G\delta_N$ cancel in the coefficient of $\ell^2$, which is
$(-\frac72+\frac74+\frac{67}4)\mathcal V(G)=15\mathcal V(G)$. So $\sum_{s\in D}\cdots\equiv15\ell^2\mathcal V(G)\pmod{\ell^3}$. The constant factor $\langle\kappa_1,\kappa_2\rangle\equiv1\pmod\ell$ does not matter, because the sum is $0$ modulo $\ell^2$. Hence
$T\equiv\ell K_0'\cdot15\ell^2\mathcal V(G)\pmod{\ell^4}$. Finally $K_0\equiv K_0'\pmod\ell$, because Wilson's theorem gives $(2n)!(\alpha_0-1)!\equiv-1$. \qed

\begin{lemma}\label{lem:Wt}
Let $\delta_N^{(j)}(x)=\sum_{i<\alpha_0}(x-i)^{-j}$ and $\delta_A^{(j)}(x)=\sum_{i\le A_2}(x-i)^{-j}$. The functions listed in Step 4 are polynomials in $\Z_{(\ell)}[x]$ of degree at most
$\ell-4$ that vanish at every $x\in[0,\alpha_0)$.
\end{lemma}

\begin{proof}
$\Omega$ has a zero of order $3$ at each $i<\alpha_0$ and each factor has a pole of order at most $2$ there. The degree of $G$ is at most $\ell-4$ and the others are smaller.
\end{proof}

Lemma~\ref{lem:G3} makes $\overline N_{\mathrm{half}}$ explicit. For $\varphi=t^{j+j'}$ we have $\psi=(-x)^{j+j'}$, so up to the unit $b$ and the unit scaling of the basis
\[
 \det\overline N_{\mathrm{half}}=(15K_0)^g\det\bigl[\mathcal V(\Omega\,x^{j+j'})\bigr]_{j,j'<g},
\]
because $\det[(-1)^{j+j'}h_{j+j'}]=\det[h_{j+j'}]$. The following closed form shows that this is a unit.

\begin{lemma}[the Hankel determinant]\label{lem:W}
Let $\gamma\ge0$, $\alpha_0=2\gamma+1$ and $g\ge1$. Then
\[
 \det\bigl[\mathcal V(D_{\alpha_0}(x)^3x^{i+j})\bigr]_{0\le i,j<g}=(-1)^{\gamma g+\lfloor g/2\rfloor}\prod_{k<\lceil g/2\rceil}h^+_k\prod_{k<\lfloor g/2\rfloor}h^-_k,
\]
where
\begin{align*}
 h^+_k&=\frac{k!\,(\gamma+k)!^3\,(2\gamma+1+k)!^3\,(k+3\gamma+1)!}{(2k+3\gamma+1)!\,(2k+3\gamma+2)!},\\
 h^-_k&=\frac{k!\,(\gamma+1+k)!^3\,(2\gamma+1+k)!^3\,(k+3\gamma+2)!}{(2k+3\gamma+2)!\,(2k+3\gamma+3)!}.
\end{align*}
All factorial arguments are at most $g+3\gamma+1$. So the determinant is nonzero, and no prime $\ell\ge g+3\gamma+2$ divides its numerator or its denominator.
\end{lemma}

By \eqref{eq:band}, $\ell=6\gamma+2g+5>g+3\gamma+1$. Since $\ell\ge19$, the factor $15K_0$ is a unit as well. So
$\ell\nmid\det\overline N_{\mathrm{half}}$. The numbers $h^\pm_k$ are the squared norms of Wilson polynomials with
parameters $(0,\gamma+1,\gamma+1,\gamma+1)$ and $(1,\gamma+1,\gamma+1,\gamma+1)$ \cite{Wilson80,KLS}. The
identification of the weight is a remark and not a step of the proof. We prove the lemma by finite algebra. Put $\varsigma=\gamma+1$ and $(y)_k=y(y+1)\cdots(y+k-1)$.

\emph{Two facts on $\mathcal V$.} (L0) $\mathcal V(f(x+1)-f(x))=f''(0)$. By linearity take $f=x^m$. Then
$\mathcal V((x+1)^m-x^m)=\sum_{e<m}\binom me eB_{e-1}=m\sum_{e'<m-1}\binom{m-1}{e'}B_{e'}$, which is $m[m=2]=(x^m)''(0)$ by $\sum_{k<N}\binom NkB_k=[N=1]$.
(R0) $\mathcal V(g(-x))=-\mathcal V(g)-g''(0)$, since $B_{e-1}=0$ for even $e\ge4$. By induction, $\mathcal V(g(x+m))=\mathcal V(g)+\sum_{j<m}g''(j)$ and
$\mathcal V(g(m-x))=-\mathcal V(g)-\sum_{j\le m}g''(j)$.

\emph{The centred functional.} Put $\mathcal Z[g]=(-1)^\gamma\mathcal V\bigl(D_{\alpha_0}(x)^3g(x-\gamma)\bigr)$ on $\Q[z]$, so that $\mathcal V(D_{\alpha_0}^3x^e)=(-1)^\gamma\mathcal Z[(z+\gamma)^e]$.
\begin{itemize}
\item \emph{Evenness:} $\mathcal Z[g(-z)]=\mathcal Z[g(z)]$. Indeed $D_{\alpha_0}(2\gamma-x)=-D_{\alpha_0}(x)$, and the reflection with $N=2\gamma$ applies, because the second derivatives of $D_{\alpha_0}^3g(x-\gamma)$ vanish at $0,\dots,2\gamma$.
\item \emph{Pearson equation:} $\mathcal Z[(z+\varsigma)^3g(z+1)]=\mathcal Z[(z-\varsigma)^3g(z)]$. Put $f=D_{\alpha_0+1}^3g(x-\gamma)$. Then $f(x+1)-f(x)=D_{\alpha_0}^3[(x+1)^3g(x+1-\gamma)-(x-\alpha_0)^3g(x-\gamma)]$, and $f''(0)=0$, so (L0) applies.
\end{itemize}
Put $F_j=(z)_j(-z)_j$, $E_k=(\varsigma+z)_k(\varsigma-z)_k$ and $F^1_j=(1+z)_j(1-z)_j$. These are even, of degree $j,k,j$ in $z^2$, and $z^2F^1_j=-F_{j+1}$.

\emph{Beta evaluation.} For $j,k\ge0$,
\[
 \mathcal Z[F_jE_k]=\mathcal Z[1]\,\frac{(\varsigma)_{j+k}(\varsigma)_j^2(2\varsigma)_k^2}{(3\varsigma)_{j+k}} .
\]
Let $g(z)=(z)_j(1-z)_j(\varsigma+z)_k(\varsigma+1-z)_k$ and $T(z)=(z+\varsigma)^3g(z+1)$. The Pearson equation gives $\mathcal Z[T]=\mathcal Z[(z-\varsigma)^3g]=\mathcal Z[-T(-z)]=-\mathcal Z[T]$ by evenness, so $\mathcal Z[T]=0$.
Evenness also gives $\mathcal Z[T]=\frac12\mathcal Z[T(z)+T(-z)]$. With $zT(z)=P(z)F_jE_k$, $P=(\varsigma+z)^2(\varsigma+k+z)(z+j)$, we get $(P(z)-P(-z))/z=2[(3\varsigma+j+k)z^2+2\varsigma j(\varsigma+k)+\varsigma^2(\varsigma+j+k)]$, hence
$\mathcal Z[F_jE_k((3\varsigma+j+k)z^2+2\varsigma j(\varsigma+k)+\varsigma^2(\varsigma+j+k))]=0$. Using $z^2E_k=(\varsigma+k)^2E_k-E_{k+1}$ and $z^2F_j=j^2F_j-F_{j+1}$ gives the two recursions
\begin{align*}
 (3\varsigma+j+k)\mathcal Z[F_jE_{k+1}]&=(\varsigma+j+k)(2\varsigma+k)^2\mathcal Z[F_jE_k],\\
 (3\varsigma+j)\mathcal Z[F_{j+1}]&=(\varsigma+j)^3\mathcal Z[F_j],
\end{align*}
which the right side satisfies. Induction on $j$ and then on $k$ gives the formula.

\emph{The mass.} $\mathcal Z[1]=(2\gamma+1)!^3\gamma!^3/(3\gamma+2)!$. Write $\mathcal Z_\gamma$ for $\mathcal Z$
to show the dependence on $\gamma$. Then $D_{2\gamma+1}(x)=\Phi_\gamma(x-\gamma)$ with $\Phi_\gamma(z)=z\prod_{i=1}^\gamma(z^2-i^2)$, and (L0) gives
$\mathcal Z_\gamma[g]=-\mathcal Z_{\gamma-1}[(z^2-\gamma^2)^3g]$. So $\mathcal Z_\gamma[1]=\mathcal Z_{\gamma-1}[u^3]$ with $u=\gamma^2-z^2=E_1$ at parameter $\gamma-1$, where $\varsigma=\gamma$. Write $u^3=E_3-(6\gamma+5)E_2+(2\gamma+1)^2E_1$ and use the Beta evaluation at $j=0$. A short computation gives
$\mathcal Z_\gamma[1]/\mathcal Z_{\gamma-1}[1]=((2\gamma+1)2\gamma)^3\gamma^3/(3\gamma(3\gamma+1)(3\gamma+2))$, and $\mathcal Z_0[1]=\mathcal V(x^3)=3B_2=\frac12$.

\emph{Splitting.} The matrix $[\mathcal Z[(z+\gamma)^{i+j}]]$ equals $U[\mathcal Z[z^{i+j}]]U^\top$ with $U$ lower unitriangular, so the determinant in the lemma is $(-1)^{\gamma g}\det[\mathcal Z[z^{i+j}]]$. By evenness $\mathcal Z[z^{\rm odd}]=0$, and ordering the indices as (even, odd) makes the matrix block diagonal. The blocks are the Hankel matrices of
$\mathcal Z_+(P)=\mathcal Z[P(z^2)]$ of size $\lceil g/2\rceil$ and $\mathcal Z_-(P)=\mathcal Z[z^2P(z^2)]$ of size $\lfloor g/2\rfloor$. Let $\mathcal K$ be a linear functional on $\Q[w]$. If $\mathcal P_m$ are monic of degree $m$ with $\mathcal K(\mathcal P_mq)=0$ for $\deg q<m$, then $\det[\mathcal K(w^{i+j})]_{i,j<r}=\prod_{m<r}\mathcal K(\mathcal P_mw^m)$. Indeed $[\mathcal K(\mathcal P_i\mathcal P_j)]=V[\mathcal K(w^{i+j})]V^\top$ with $V$ unitriangular.

\emph{The orthogonal polynomials.} For $m\ge0$ put
\[
 \mathcal Q_m=\sum_{k\le m}\frac{(-m)_k(m+3\varsigma-1)_k}{k!(\varsigma)_k(2\varsigma)_k^2}E_k,\qquad \mathcal R_m=\sum_{k\le m}\frac{(-m)_k(m+3\varsigma)_k}{k!(\varsigma+1)_k(2\varsigma)_k^2}E_k .
\]
For $j\le m$ the Beta evaluation gives
\[
 \mathcal Z[\mathcal Q_mF_j]=\mathcal Z[1]\frac{(\varsigma)_j^3}{(3\varsigma)_j}\sum_{k}\frac{(-m)_k(m+3\varsigma-1)_k(\varsigma+j)_k}{k!(\varsigma)_k(3\varsigma+j)_k}.
\]
This sum is balanced. Pfaff--Saalschütz \cite[Theorem~2.2.6]{AAR99} with $\alpha=m+3\varsigma-1$, $\beta=\varsigma+j$ and $\gamma'=\varsigma$ gives a product containing $(-j)_m$, which is zero for $j<m$. So $\mathcal Q_m$ is orthogonal for $\mathcal Z_+$ to all polynomials in $z^2$ of degree less than $m$. Let $\hat{\mathcal Q}_m$ be the monic multiple of $\mathcal Q_m$ in the variable $w=z^2$. At $j=m$ we divide by the
leading coefficient $(m+3\varsigma-1)_m/((\varsigma)_m(2\varsigma)_m^2)$ of $\mathcal Q_m$ and by the leading
coefficient $(-1)^m$ of $F_m$. The product then evaluates to $\mathcal Z_+(\hat{\mathcal Q}_mw^m)=h^+_m$. For the odd block, $\mathcal Z_-(\mathcal R_mF^1_j)=-\mathcal Z[\mathcal R_mF_{j+1}]$. The same argument with $\gamma'=\varsigma+1$ and $\beta=\varsigma+j+1$ gives orthogonality for $j<m$ and $\mathcal Z_-(\hat{\mathcal R}_mw^m)=-h^-_m$. The factorial arithmetic uses
$(3\varsigma)_m=(m+3\gamma+2)!/(3\gamma+2)!$ and $(m+3\varsigma-1)_m=(2m+3\gamma+1)!/(m+3\gamma+1)!$. The product of the norms, with the sign $(-1)^{\gamma g}$ from the shift and $(-1)^{\lfloor g/2\rfloor}$ from the odd block, is the formula. The largest factorial arguments are
$2m+3\gamma+2\le g+3\gamma+1$ in the even block and $2m+3\gamma+3\le g+3\gamma+1$ in the odd block. \qed

\emph{Remark.} Put $W(y)=|\Gamma(\gamma+1+iy)|^6\cosh(\pi y)/\pi$. A Barnes-integral evaluation gives $\mathcal Z[g]=\int_{\R}g(iy)W(y)\,dy$, and $W$ is the Wilson weight with parameters $(0,\gamma+1,\gamma+1,\gamma+1)$. Its odd part $y^2W$ is the Wilson weight with $(1,\gamma+1,\gamma+1,\gamma+1)$. This explains the two families $h^+,h^-$. We do not use it.

\subsubsection{Proof of Theorem~\ref{thm:B83}}
Lemma~\ref{lem:blockA}, Lemma~\ref{lem:blockC} and the cross entries give $\det\overline N_{\mathrm{int}}\ne0$ for
$\ell\notin\{2,13,24359\}$. Lemmas~\ref{lem:G3} and~\ref{lem:W} give $\det\overline N_{\mathrm{half}}\ne0$ for
$\ell\ge19$. So $\det\overline N(0,0)\ne0$. By \eqref{eq:target} and Lemma~\ref{lem:PD}, the valuation of
$J_n(a,b,X)$ is $\mathrm{w}_{n,\ell}$ at every point of the theorem. In particular $J_n(a,b,X)\ne0$. \qed

%% =====================================================================
\section{Proof of Theorem~\ref{thm:main}}\label{sec:proof}

Put $D=33.4401$, $\Lambda=28.68723663$, $\eta=D-\Lambda=4.75286337$, $\rho=42.94856306$ and $\kappa=3$. The modulus is
$M=126$, and there are four tied families:
\begin{itemize}
\item the band: all unit classes, $c=\frac83$, $c'=3$, $R(\ell)=\{n:8n<3\ell,\ \ell+1<3n\}$;
\item Lai's primes: $S=\{s:s\equiv1\pmod{18}\}$, $c=c'=\frac92$, $\ell=\frac92n+1$;
\item the first line: $S=\{s:s\equiv5\pmod7\}$, $c=c'=\frac73$, $\ell=\frac73n+\frac13$;
\item the second line: $S=\{s:s\equiv3\pmod7\}$, $c=c'=\frac73$, $\ell=\frac73n+\frac23$.
\end{itemize}
The band has the required property with $C=2$ for $\ell\ge100$. Indeed the interval $(\frac{\ell+1}3,\frac{3\ell}8)$
has length at least $3$, so it contains an integer within $2$ of every $y\in[\frac\ell3,\frac{3\ell}8]$. Each of the
three other sets $S$ has $6$ of the $36$ unit classes. Lai's set meets each of the two lines in one class, and the
two lines are disjoint. The finite exceptional set $E$ consists of the primes below $100$ and the primes $1039$ and
$24359$.

\emph{Part (a).} Since $a\zeta(3)+b\pi^4=a\zeta(3)+90b\zeta(4)$, it suffices to prove the bound of part (a) for
$a\zeta(3)+b\zeta(4)$ with $a,b\in\Z$ and $b\ne0$. Let $\xi=a\zeta(3)+b\zeta(4)$ and $P_n=Q_n$. By \eqref{eq:scaled} and
Theorem~\ref{thm:content}, $q^{3n}m_nQ_n(p/q)=m_nJ_n(qa,qb,p)\in\Z$, so (a) of Theorem~\ref{thm:criterion} holds with
$d_n=m_n$. Corollary~\ref{cor:families} gives (b). Take $W(p,q)=q$ and $K=1$. For each family add the prime divisors of
$b$ to the exceptional set, and let $\mathrm{Bad}(\ell,p,q)$ be $\ell^4\mid q$, with $e=4$. Corollaries~\ref{cor:B83},
\ref{cor:NV} and~\ref{cor:NV7} give (c). Put $u_0=3/\eta$ and $\beta=7/800$. The windows $(\frac83u,3w]$, $(\frac92u,\frac92w]$ and
$(\frac73u,\frac73w]$ with $w=u_0+\beta$ are pairwise disjoint, because $\frac73w<\frac83u_0$ and $3w<\frac92u_0$. So
\begin{multline*}
 \cM(u_0,u_0+\beta)=4\Bigl[\Bigl(\frac{u_0}3+3\beta\Bigr)+\frac16\cdot\frac92\beta+2\cdot\frac16\cdot\frac73\beta\Bigr]\\
 =\frac4\eta+\frac{163}9\beta=1.00007\ldots>1 .
\end{multline*}
Theorem~\ref{thm:criterion} applies, and $\rho(3/\eta+\beta)=27.48487\ldots<27.49$.

\emph{Part (b).} Let $\xi=y^*=\zeta(3)/\zeta(4)$ and $P_n=f_n$. As above,
$q^{3n}m_nf_n(p/q)=m_nJ_n(q,-p,0)\in\Z$, and Corollary~\ref{cor:families} gives (b). Since $1<y^*<2$, every $p/q$
with $|y^*-p/q|<1$ has $0<p<3q$, and we take $W(p,q)=|p|$ and $K=3$. Let $\mathrm{Bad}(\ell,p,q)$ be $\ell\mid p$, with
$e=1$. Corollaries~\ref{cor:B83}, \ref{cor:NV} and~\ref{cor:NV7} give (c). With $\beta=9593/50000$ the windows of the band and of the
two lines overlap, and the other windows are disjoint from them. The band meets each line in a window of length $\frac73w-\frac83u_0=\frac13(7\beta-u_0)$ and
density $\frac16$, so
\begin{multline*}
 \cM(u_0,u_0+\beta)=\Bigl(\frac{u_0}3+3\beta\Bigr)+\frac16\cdot\frac92\beta+2\cdot\frac16\cdot\frac73\beta-2\cdot\frac16\cdot\frac13(7\beta-u_0)\\
 =\frac{4u_0}9+\frac{15}4\beta=1.0000076\ldots>1 .
\end{multline*}
Since $\rho(3/\eta+\beta)=35.34918\ldots$, Theorem~\ref{thm:criterion} with $\mu=35.3495$ gives
$|y^*-p/q|>q^{-35.3495}$ for all integers $p$ and all large $q$.
Finally $\zeta(3)/\pi^4=y^*/90$, because $\zeta(4)=\pi^4/90$. Lemma~\ref{lem:scale} with $r=1/90$, $\mu'=35.3495$
and $\mu=35.35$ gives part (b) of Theorem~\ref{thm:main}. \qed

%% =====================================================================
\section{Remarks}\label{sec:remarks}

The numbers in this section are numerical evidence. They are not part of the formal proof.

\emph{The decay.} The cell computation of Theorem~\ref{thm:decay} gives $D\ge33.44015$. Optimizing the measure on wedge
contours suggests that the best value of this method for Lai's determinants is close to $33.46$. On a vertical line,
weak duality caps the method at about $33.25$. Bending the contour is therefore necessary for the constant of
Theorem~\ref{thm:decay}. Lai's exponent $31.61$ holds at the constant only, and Theorem~\ref{thm:decay} improves it
on a neighbourhood.

\emph{The denominators.} At more than $3000$ pairs $(n,p)$ with $n\le200$ we computed $v_p(J_n)$ exactly. It never lies
below the bounds of Section~\ref{sec:local}, and in the range $\sqrt{8n}<p<3n$ it equals $\Lex(n,p)+x_1+x_2$ except at a
few primes where a class constant of the proofs vanishes modulo $p$. This suggests that $\Lambda=28.6872366\ldots$ is
the true exponential rate of the denominators of $J_n(a,b,X)$ as a polynomial in three variables.

\emph{Nonvanishing.} The tolerance $\ell^4\nmid q$ is best possible for all three families: at the tested pairs, the roots of
$J_n(1,1,X)$ have $v_\ell(X)=-4$. For primes $5n/2<\ell<8n/3$ the reduced half block is the Hankel matrix of a weight on two
intervals of residues. We know no closed form for its determinants, which have large prime factors. At $124$ pairs
$(n,\ell)$ with $n<90$ the prime $\ell$ does not divide it. For one size $n$ the primes of the band carry the mass
$\sum_\ell e\log\ell\sim\frac43n$ for the plane and $\frac13n$ for the ratio. A mass of $1.585n$ would give the
exponent $27.11$. The exponent of part (b) is larger than that of part (a) because the coefficient $b=-p$ of the
ratio family grows with $q$. So only one power of each tied prime is available.

\emph{The axes $a=0$ and $b=0$.} For $b=0$ the number is $a\zeta(3)$. Theorem~\ref{thm:main} does not apply, and the
known measures for $\zeta(3)$ are much smaller. For $a=0$ the number is a rational multiple of $\pi^4$, and the bound
$12.51$ of Marcovecchio and Zudilin is better than ours.

%% =====================================================================
\section{Formalization}\label{sec:lean}
\raggedbottom

Theorem~\ref{thm:main} is formalized in Lean~4 \cite{Lean4} with Mathlib \cite{Mathlib}, together with every theorem,
proposition and lemma its proof uses. The remarks, Theorem~\ref{thm:C01}, the exact valuations stated after
Theorems~\ref{thm:NV}, \ref{thm:NV7} and~\ref{thm:B83}, and the remark on Wilson weights after Lemma~\ref{lem:W} are
not part of the formalization. The development is in the directories
\lean{CatalanCore/Lai} and \lean{CatalanCore/Measure} of the repository
\begin{center}
\FormalizationLink
\end{center}
Theorem~\ref{thm:main} consists of the theorems
\begin{center}
{\small\texttt{LaiIrrationalityMeasure.lai\_plane\_measure\_27\_49}}\\
{\small\texttt{LaiIrrationalityMeasure.lai\_ratio\_measure\_35\_35}}
\end{center}
in \lean{SolutionLai.lean}, and Corollary~\ref{cor:mobius} is
\begin{center}
{\small\texttt{LaiIrrationalityMeasure.lai\_ratio\_mobius\_measure\_35\_35}.}
\end{center}
Their statements use only Mathlib. The numbers $\zeta(3)$ and $\zeta(4)$ are written as the
series $\sum_{k\ge0}(k+1)^{-3}$ and $\sum_{k\ge0}(k+1)^{-4}$, and $\pi$ is Mathlib's \lean{Real.pi}. The first theorem
states, for integers $a$ and $b\ne0$, that $q^{-27.49}<|a\zeta(3)+b\zeta(4)-p/q|$ for all integers $p$ and all
sufficiently large $q$. Applied to $(a,90b)$ it gives part (a), since $\pi^4=90\zeta(4)$. This is the only place
where the formal statement differs in form from Theorem~\ref{thm:main}. The second states $q^{-35.35}<|\zeta(3)/\pi^4-p/q|$ in the same form. The modules contain no
\lean{sorry} and no \lean{native_decide}. The command \lean{\#print axioms} reports only \lean{propext},
\lean{Classical.choice} and \lean{Quot.sound}. The proof is kernel-checked. It has not been run through the
independent checker \lean{lake comparator} \cite{Comparator}.

The prime number theorem in arithmetic progressions is derived from the Wiener--Ikehara theorem of the PNT+ project
\cite{PNTPlus} and Mathlib's results on primes in progressions. The proofs were checked with Lean
\texttt{v4.35.0-rc2} and the matching Mathlib.

The formal proof differs from the text in one point of organization. It applies the criterion to the families
$m_nQ_n$ and $m_nf_n$, which have integer values at $p/q$ after multiplication by $q^{3n}$, with content exponent $0$
and decay exponent $D-\Lambda$. This is equivalent to the form of Theorem~\ref{thm:criterion} used in
Section~\ref{sec:proof}.

The results of the paper correspond to the following declarations. Their files are under \lean{CatalanCore}. We give
the names with a short prefix only. The full namespaces are \lean{CatalanCore.Lai}, \lean{CatalanCore.Measure},
\lean{CatalanCore.Content} and \lean{CatalanCore.NT}.
\begin{itemize}
\item Proposition~\ref{prop:finite} and \eqref{eq:homog}: \lean{laiR_partial_fractions}, \lean{laiL_eval_moment} and
 \lean{isHomogeneous_laiJ}. The identity $\zeta(4)=\pi^4/90$ is \lean{zeta4_eq_pi}.
\item Theorem~\ref{thm:criterion}: \lean{measure_of_family_tiedNV_index}, with the structure \lean{TiedNV}, whose families have a slope interval $[c,c']$, the mass
 condition \lean{NT.UnionWindowOK} and the prime number theorem in progressions \lean{NT.pnt_class} and
 \lean{NT.theta_interval_asymp}. Lemma~\ref{lem:scale}: \lean{plain_measure_smul}. Lemma~\ref{lem:mobius}:
 \lean{plain_measure_all_den} and \lean{plain_measure_mobius}. Corollary~\ref{cor:mobius}: \lean{lai_mobius_of}.
\item Theorem~\ref{thm:dual}: \lean{lai_decay_dual}. Theorem~\ref{thm:decay} and Corollary~\ref{cor:families}:
 \lean{lai_wedge_D}, \lean{lai_plane_decay} and \lean{lai_ratio_decay}, applied to the certificate \lean{laiDecayData}.
\item Proposition~\ref{prop:crude}: \lean{lai_vp_lower}, \lean{lai_vp_nonneg} and \lean{lai_small_primes}.
 Theorem~\ref{thm:O}: \lean{thmO}. Theorem~\ref{thm:B7}: \lean{band7_proof}.
\item Lemma~\ref{lem:S}: \lean{lemmaS} and \lean{lemmaT}. Theorem~\ref{thm:L1}: \lean{lai_L1}.
 Theorem~\ref{thm:L2}: \lean{lai_L2}.
\item Proposition~\ref{prop:bridge}: \lean{bridge_proof}. Lemma~\ref{lem:SM}: \lean{c01_partial_le}.
 Lemma~\ref{lem:bands} and Theorem~\ref{thm:content}: \lean{sum_signed_le} and \lean{lai_content}.
\item Theorem~\ref{thm:NV} and Corollary~\ref{cor:NV}: \lean{lai_nv_core}, \lean{lai_nv_plane} and \lean{lai_nv_ratio}.
 Theorem~\ref{thm:NV7} and Corollary~\ref{cor:NV7}: \lean{nv7_core}, \lean{nv7_plane} and \lean{nv7_ratio}.
\item Theorem~\ref{thm:B83}: \lean{B83Main.b83_nonvanishing}, with the corollaries \lean{nv_plane_B83} and \lean{nv_ratio_B83} for Corollary~\ref{cor:B83}.
 Lemma~\ref{lem:U83}: \lean{u83_A} and \lean{u83_C}. Lemmas~\ref{lem:blockA} and~\ref{lem:blockC}: \lean{det_MA} and \lean{det_MC}, with \lean{det_intBlock_band_ne}.
 Lemma~\ref{lem:entries}: \lean{lb06_point} and the weighted bounds in the file \lean{CatalanCore/Lai/B83Weighted.lean}. Lemma~\ref{lem:G3}: \lean{lemma_G3_monomial}, with Wolstenholme's theorem
 \lean{wolstenholme_one} and \lean{wolstenholme_two}. Lemma~\ref{lem:W}: \lean{wilson_hankel}, \lean{wilson_hankel_padic_unit} and the
 Pfaff--Saalschütz summation \lean{saalschutz}. The unit of the half block: \lean{half_block_unit}. The assembly of the Gram determinant: \lean{det_gramF_ne_zero}.
\item Section~\ref{sec:proof}: \lean{lai_roots_of}.
\end{itemize}
The certificate data are the coefficients of $H$, the $40$ cells, the trigonometric values at the cell centres, and
the $22977$ terms of the partial sum in Lemma~\ref{lem:SM}. They were computed with floating-point optimization and
the interval arithmetic of Arb \cite{Arb}, and Lean re-checks them with rational arithmetic. They play no role in the
proof beyond satisfying the checks.

\section*{Disclosure of AI use}
The mathematics, the Lean formalization, the certificate generators and the text of this paper were produced by AI
systems at the request of the author. The author posed the problems, directed the work and is responsible for the
paper. The systems were Claude Opus~5.5, Claude Sonnet~5.5 and Claude Haiku~5.5 (Anthropic), running in Claude Code,
and GPT-6.1 Sol (OpenAI), running in Codex. The construction of the determinants and the linear independence of $1$,
$\zeta(3)$ and $\pi^4$ are due to Lai \cite{Lai2026}.

\subsection*{Acknowledgements}
We thank the authors of Mathlib and of the PNT+ project.

\bibliographystyle{amsplain}
\bibliography{refs}

\end{document}
